\documentclass[11pt]{article}
\usepackage[margin=2.5cm]{geometry}
\usepackage{amsmath,amssymb}
\usepackage{booktabs,longtable,array}
\usepackage[utf8]{inputenc}
\usepackage[semicolon]{natbib}
\usepackage{enumitem}
\usepackage[hidelinks]{hyperref}
\setlist{itemsep=2pt,topsep=3pt}
\setlength{\parskip}{4pt}
\setlength{\emergencystretch}{3em}
\newcommand{\code}[1]{\texttt{#1}}
\newcommand{\E}{\mathbb{E}}
\newcommand{\1}{\mathbf{1}}
\newcolumntype{L}[1]{>{\raggedright\arraybackslash}p{#1}}

\title{A Monthly Two-Earner Model of Joint Job Loss\\[4pt]
\large Model document for Paper A (\emph{The Insurer Is Hit Too}), with the PI's decisions}
\author{Paper A --- model documentation}
\date{Design state of 18 October 2026, brought up to the code of 19 October 2026 (\code{IGE/paperA\_model/})}

\begin{document}
\maketitle

\begin{abstract}
\noindent This document states the structural model of Paper A, in the order in which the simulator runs it, so that
the state space, the timing, the policy schedules, the parameters and the moments can be checked against the code
(\code{IGE/paperA\_model/}; where the two differ, the code is right and this document is corrected). Part~I sets out the model as an economics paper would: the environment, preferences, the budget
constraint and the government, the labor market and its shocks, the marriage, the timing within the month, the
household's problem as Bellman equations, the children block and the solution method. Part~II sets out the
estimation: every parameter with its status (calibrated, with the source, or estimated, with the moments that identify
it), the targets with their sample windows, the untargeted validation, the protocol for running the estimator of
\citet{blundell2016consumption} (BPS) on simulated households with its pass criteria, the five counterfactuals, and the
PI's decisions of 18 October 2026 and the six taken on 18--19 October after the first evaluations of the coded model
(Section~\ref{sec:decisions}; none is left open). Every modelling choice that rests on
a data moment carries a footnote with the number and the script that produced it. The document contains no model results: the numbers in it are data moments of Paper A or
fixed inputs.
\end{abstract}

\tableofcontents
\newpage

% #####################################################################################################
\part{The model}

% =====================================================================================================
\section{Environment}\label{sec:env}

\paragraph{Households and time.} The economy is populated by married couples. Time is discrete and the period is a
month.\footnote{In the CPS (1994--2023) the wife's excess risk of losing her job is 4.0 percentage points
(0.040$^{***}$, s.e.\ 0.002, $p<0.001$) in the month her husband loses his, against 0.6--0.8 points in each of the
three months before and after (all $p<0.001$): the common shock is a same-month coincidence. Wives who were not working
search at once (unemployed and looking $+8.8$ points, 0.088$^{***}$, $p<0.001$, the month after) and their employment
builds up over about six months (SIPP: 0.019$^{***}$ in the month of the husband's separation, 0.054$^{***}$ after six
months, $p<0.01$). Unemployment insurance lasts 26 weeks. In an annual period the common and the permanent components
of the risk collapse into one same-year correlation (PSID: she is displaced in 10.2\% of the years he is, against 3.9\%),
and the CPS leads and lags cannot be used; averaged over a year, the joint losses are about 2\% of the displaced family's
income loss. Most US earnings risk is month-to-month \citep{ganong2026earnings}. Code: \code{IGE/30.CPSJointJobLoss.py},
\code{IGE/37.CPSJointLossChecks.py}, \code{IGE/40.SIPPJointLoss.py}, \code{IGE/25.JointJobLoss.py}.}
A couple enters the model in January of the first calendar year in which the corresponding PSID couple is observed
(Section~\ref{sec:init}) and is followed until the wife is 60 (Decision~\ref{od:period}): with the wife aged 25 at the
earliest, $t = 0, \dots, T$ indexes months, $T = 419$, the wife's age is $25 + \lfloor t/12 \rfloor$, and $t$ is a
multiple of 12 in every January, so model years are calendar years. Age-dependent inputs and the transitory wage shocks
change in January. All money is in 2000 dollars a month, the unit of Paper A's PSID tables.

\paragraph{Decision maker.} The couple is a unitary household: it maximises one utility function, in which the spouses
share consumption and each spouse's work enters as a disutility. There is no bargaining. The joint-loss moments carry no
information on a Pareto weight, and the policy questions are about the household's resources in the months of the
shock; a collective model would add a parameter the data cannot separate. After a separation the wife is the only
decision maker of her household; the husband leaves the model.

\paragraph{Heterogeneity.} Couples differ in four permanent characteristics, all known to the couple:
\begin{itemize}
  \item race $r \in \{B, W\}$, which enters only through the environment (displacement hazards, the search friction, the
  separation hazard, the initial conditions and the inputs), never through preferences;\footnote{The joint loss is the
  same in both races (CPS excess Black 0.044$^{***}$ on a base of 0.60\%, White 0.038$^{***}$ on 0.54\%; difference 0.006,
  $p=0.309$; Displaced Worker Supplement difference $-0.001$, $p=0.894$), while the environment differs: husbands'
  annual displacement rate is 8.0\% (Black) against 4.8\% (White), and the separation response to his job loss is
  $-4.5$ against $-1.9$ points. Common preferences make the racial differences in the wife's response a test of the
  environment. Code: \code{IGE/30.CPSJointJobLoss.py}, \code{IGE/39.CPSDisplacedWorkers.py},
  \code{IGE/19.FatherDisplacement.py}. Decision~\ref{od:prefs}.}
  \item workplace match $s \in \{0, 1, 2\}$, which sets the common component of displacement risk: $s=2$ if the spouses
  share a detailed industry, a detailed occupation and a class of worker (the closest the public data come to a shared
  employer), $s=1$ if they share a detailed industry but not that cell, and $s=0$ otherwise;\footnote{In the CPS
  (1994--2023, all job losses, state $\times$ month fixed effects) the wife's same-month excess is 0.337$^{***}$ (s.e.\
  0.015) in the shared-employer cell, which holds 6.3\% of couples and 35\% of the raw excess; 0.211$^{***}$ (s.e.\
  0.009) among couples in the same detailed industry (17.7\% of couples, 53\% of the excess; same major industry 57\%);
  and 0.026$^{***}$ (s.e.\ 0.001) among all couples outside the shared-employer cell. If 10--12\% of dual-earner couples
  work at the same establishment \citep{hyatt2019coworking}, a shared employer carries a third to a half of the excess,
  and the same industry more than half. In the PSID, where the spouses work in the same one of 16 sectors (15\% of
  couples) the husband's displacement raises the wife's risk of displacement by a further 20 points (0.203$^{***}$, s.e.\
  0.055), in the same three-digit industry (8\%) by 28 points (0.282$^{***}$); in the SIPP the excess rises by 22 points
  (0.220$^{***}$) in the same detailed industry. The PSID couple's $s$ at entry is coded with the CPS definition from both
  spouses' three-digit industry, occupation and class of worker (\code{IGE/60.PaperATargets.py}). Code:
  \code{IGE/37.CPSJointLossChecks.py}, \code{IGE/25.JointJobLoss.py}, \code{IGE/40.SIPPJointLoss.py}.}
  \item a couple risk type $\tau \in \{L, H\}$, which scales both spouses' own displacement hazards and stands for the
  assortative part of the risk (couples in which both spouses hold riskier jobs);\footnote{Wives of men who will be
  displaced four to six years later are already more often displaced in the placebo window (0.033$^{**}$, $p=0.012$);
  stacking both windows, the permanent couple component of the excess is 2.0 points (0.020$^{**}$, $p=0.012$) and the
  contemporaneous component 3.3 points (0.033$^{**}$, $p=0.036$). In the CPS the elevation in the months around his job
  loss (0.6--0.8 points) is flat, as a permanent type predicts. Code: \code{IGE/25.JointJobLoss.py},
  \code{IGE/38.PSIDJointLossChecks.py}, \code{IGE/30.CPSJointJobLoss.py}.}
  \item marital status $m \in \{M, S\}$ (married, separated; separation is absorbing).
\end{itemize}
Education does not index the solution: it enters through each spouse's initial human capital level and initial wealth
(Section~\ref{sec:init}); an education dimension is added only if the fit by education fails (Decision~\ref{od:educ}).

\paragraph{Choices.} Each month the household chooses consumption and saving subject to a borrowing limit, the wife's
hours $\ell_w \in \{0, P, F\}$ when she has a job (zero means she leaves it), the husband's hours $\ell_h \in \{P, F\}$
when he has a job, and whether to take up food stamps (SNAP) when eligible. Unemployment insurance is collected
whenever a spouse is entitled to it. There is no search-effort choice and no separate search-effort state
(Decision~\ref{od:search}): a non-employed spouse finds a job with a probability that measures the search friction, and
a wife who finds one decides whether to work.\footnote{``Looking for
work'' is therefore a measurement of the model, not a choice: a non-employed wife is classified as looking when she
would accept a job in her current state with probability above one half. The monthly CPS jump in looking ($+8.8$
points the month after his loss) is used as validation of this mapping. The PSID's annual interview-week status cannot
discipline a search choice: a year after the job loss looking is unchanged ($-0.001$, s.e.\ 0.008). Code:
\code{IGE/30.CPSJointJobLoss.py}, \code{IGE/19.FatherDisplacement.py}.}
The husband works whenever he has a job.\footnote{Married men who were employed the month before are still employed in
99.5\% of cases a month later and in 96.0\% a year later when they have no involuntary separation (SIPP, husbands
without a separation in the month); their exits outside displacement are rare and are left out. Code:
\code{IGE/40.SIPPJointLoss.py}.}

\paragraph{What is exogenous.} Children's arrival and ageing, the workplace match, the couple type and the policy
schedules are exogenous. Displacement and the separation are exogenous shocks whose probabilities depend on the state.
Employment, hours, savings, SNAP take-up and, through them, human capital, UI entitlement and consumption are
endogenous.

\paragraph{Relation to \citet{bacher2025joint}.} The closest model is the monthly life-cycle model of joint search of
Bacher, Gr\"ubener and Nord. Two spouses of the same age live 480 working months and 120 retired; each is employed,
unemployed with benefits, unemployed without benefits and searching at a cost, or out of the labor force (16 joint
states); benefits replace half of earnings up to a cap and are lost with probability $1/6$ a month; each spouse has a
human capital level on a 12-point grid that rises at work and falls out of it; the couple saves in a risk-free bond
without borrowing; discrete choices carry extreme-value taste shocks and are solved with the endogenous grid method for
discrete-continuous choice; job-finding (their ``arrival'') rates are set by firms' vacancy posting in submarkets
indexed by the household's state; the parameters are calibrated. The labor-market block here follows theirs: a monthly
period, employment and non-employment with and without unemployment insurance that expires at a constant hazard, a
human capital level for each spouse that accumulates at work and depreciates out of it, savings, and the same solution
method; their 480-month two-member solution is the precedent for the monthly period (Decision~\ref{od:period}). It
departs in three ways. (i) The spouses' displacement risk has a common component, larger for couples who share an
employer or an industry, and a permanent couple component; their separation shocks are independent across spouses
(their footnote 19), which bounds the ratio of the wife's job-loss probability in the husband's loss month to other
months at about 1.25, against about 10 in the CPS.\footnote{With independent hazards that depend only on the human
capital level, the ratio is $1 + \mathrm{cov}(\delta_w, \delta_h)/(\E\delta_w\,\E\delta_h)$; with their hazards between
1.0\% and 5.0\% a month and a mean near 2\%, perfectly correlated levels give at most about 1.25. Paper A's CPS ratio is
5.45\% against 0.54\% (\code{IGE/30.CPSJointJobLoss.py}); bound in \code{IGE/proposals/fable\_bacher2025\_paperA.md},
Section 2.3.} (ii) The search friction is exogenous and race-specific; their endogenous job-finding rates answer why
spousal insurance is available to the young and not the old, a question this model does not ask. (iii) The government
is a set of schedules (a progressive tax, the Earned Income Tax Credit, food stamps, a childcare price and a consumption
floor besides unemployment insurance), the couple has children whose ages set the cost of the wife's work (their
age-dependent disutility of non-employment, their footnote 15, stands for the same thing), the marriage can end, and
couples differ by race. All endogenous parameters are estimated jointly by the simulated method of moments with
standard errors, where theirs are calibrated. Their added-worker effect by age is an untargeted check
(Section~\ref{sec:valid}), and their unemployment-insurance duration experiment is reported next to counterfactual
(ii) (Section~\ref{sec:cf}).

% =====================================================================================================
\section{Preferences}\label{sec:prefs}

The married household's flow utility in month $t$ is
\begin{equation}\label{eq:uM}
  U^M = u\!\left(\frac{c}{s_M\, e(k)}\right) - \phi_w(\ell_w, k) - \phi_h(\ell_h) - \kappa_S\, S,
  \qquad u(x) = \frac{x^{1-\gamma} - 1}{1-\gamma},
\end{equation}
and the separated wife's is $U^S = u\big(c / e(k)\big) - \phi_w(\ell_w, k) - \kappa_S S$.

\paragraph{Consumption.} Household spending $c$ is divided by the couple's equivalence scale $s_M$ and the children's
scale $e(k)$, which adds the McClements share of the youngest child's age bucket. Risk aversion is $\gamma = 1.5$. Both
scales and $\gamma$ are calibrated and are the same as in the author's lifecycle model of women's careers, which shares
this code base.

\paragraph{The wife's work.} Her disutility of work is
\[
  \phi_w(\ell, k) = \1\{\ell = F\}\,\phi_F + \1\{\ell = P\}\,\phi_P + \1\{\ell > 0\}\,\theta(k),
  \qquad \theta(0) = 0,\; \theta(1) = \theta_{05},\; \theta(2) = \theta_{617},
\]
where $k \in \{0, 1, 2\}$ means no child under 18, a youngest child aged 0--5, and a youngest child aged 6--17. The
child term is the same for part-time and full-time work (parsimony); its dependence on the child's age is what the
data require.\footnote{After the husband's displacement, wives with a youngest child under six work 112 more hours a
year (s.e.\ 31) and are 4.1 points more likely to be employed (0.041$^{**}$, s.e.\ 0.018); with a youngest child of
school age their hours do not rise ($-26$, s.e.\ 36) and their employment falls by 3.8 points ($-0.038^{**}$, s.e.\
0.018). Among wives who kept their own job the gross response is 97 hours when the youngest child is under six
($p=0.032$) and 47 hours on average ($p=0.165$). PSID stacked event study, years 1--3. Code:
\code{IGE/19.FatherDisplacement.py}, \code{IGE/38.PSIDJointLossChecks.py}.}

\paragraph{The husband's work.} When employed he works part time or full time; $\phi_h(F) = \phi_h$ and $\phi_h(P) = 0$.
His hours choice exists so that the BPS estimator, which is identified from the intensive margin of both spouses, does
not return zero by construction (Section~\ref{sec:bps}).

\paragraph{Fixed cost of work.} Each employed spouse costs the household $F$ a month in consumption units (in the budget
constraint, Section~\ref{sec:gov}). It produces the extensive-margin complementarity of consumption and the wife's work
that BPS measure, and it means that a joint loss lowers spending both through income and through the end of two
work-related expenditure flows.\footnote{BPS's conditional Euler equations (their Table 8): consumption rises with the
wife's participation (0.388$^{**}$) and falls with her hours ($-0.241^{**}$); they attribute it to fixed costs of work
(p.~422). Spousal leisure complementarity, the other nonseparability in BPS ($\eta_{h_1,w_2} = 0.16$,
$\eta_{h_2,w_1} = 0.33$), is not modelled and is reported as a known difference.}

\paragraph{SNAP stigma.} Taking up SNAP costs $\kappa_S$ in utility, so that take-up among eligible households is
incomplete and rises after a job loss.\footnote{Food stamp receipt rises by 4.2 points for Black (0.042$^{**}$, s.e.\
0.019) and 3.3 points for White families (0.033$^{***}$, s.e.\ 0.010) in years 1--3 after the father's displacement,
while AFDC/TANF receipt barely moves (0.010 and 0.007, not significant). Code: \code{IGE/19.FatherDisplacement.py}.}

\paragraph{Taste shocks.} Each discrete alternative $d = (\ell_w, \ell_h, S)$ carries an i.i.d.\ type-I extreme value
shock with scale $\sigma_\varepsilon$, drawn each month after the transitory wage shocks and before the choice. The
scale is fixed at $\sigma_\varepsilon = 0.05$ and not estimated (Decision~\ref{od:sigma}).\footnote{A larger taste shock
does not move the wife's exit rate by the order of magnitude the data need. At the start values (one-month CPS flows,
before the extensions) $\sigma_\varepsilon = 0.05$, 0.10, 0.20 and 0.30 give a monthly exit rate of 0.26\%, 0.35\%,
0.29\% and 0.20\% against 2.6\%, while the employment of wives with no child under 18 rises from 0.95 to 0.97, 0.99 and
0.99 (data 0.81), the entry rate of White wives from 0.043 to 0.070, 0.101 and 0.113 (data 0.075), and the added-worker
response in employment falls from 0.006 to $-0.003$ (data 0.042): the option value of holding a job grows with
$\sigma_\varepsilon$, so wives keep their jobs. Objective 50{,}920, 44{,}860, 50{,}652, 57{,}938. Code:
\code{IGE/paperA\_model/results/diagnostics\_20261018/profile\_sigma\_ridge.py}; \code{results/diagnostics\_20261018.md},
A; the option \code{estimate.py --free-extra sigma\_eps} frees it.}

\paragraph{Discounting.} The monthly discount factor is fixed at $\beta = 0.96^{1/12}$ (annual 0.96); the
precautionary motive is disciplined by the estimated borrowing limit instead (Decision~\ref{od:credit}).

% =====================================================================================================
\section{The labor market and shocks}\label{sec:labor}

\subsection{Wages and human capital}

Spouse $j \in \{w, h\}$ has a human capital level $h_j \in \{1, \dots, 7\}$. Its hourly wage in month $t$ is
\begin{equation}\label{eq:wage}
  \ln w_{jt} = \mu_j(r) + \psi_j\,(h_{jt} - 4) + \nu_{jt},
\end{equation}
where $\mu_j(r)$ is a race-specific level, $\psi_j$ the log-wage step between adjacent levels and $\nu_{jt}$ a transitory
shock drawn once per calendar year (Decision~\ref{od:trans}). The wage-process parameters are calibrated to BPS's Table 3
(Decision~\ref{od:wageproc}). Monthly earnings are $y_j = w_j\, H(\ell_j)$ with $H(F) = 173$ and $H(P) = 87$ hours. There is no separate age
profile: wage growth over the life cycle comes from human capital accumulation at work, and the bounded grid flattens
it at older ages. The PSID age profile of married men's earnings, an exogenous input in the lifecycle model that shares
this code, becomes a target here (its growth) and a validation (its level).

The human capital level of each spouse changes at the end of each month in three steps, applied in this order and
truncated to $\{1, \dots, 7\}$:
\begin{enumerate}
  \item \emph{Accumulation and depreciation.} If $j$ worked this month and was not displaced, $h_j$ rises by one level
  with probability $\pi^{up}_j(\ell_j, h_j) = \pi^{up}_j\,(h_j/4)^{-\eta_{up}}$ at full time and $\kappa_{up}$ times
  that at part time (capped at one). If $j$ did not work this month, $h_j$ falls by one level with probability
  $\pi^{dn}_j$. Accumulation slows as the level rises: the level-dependent form of \citet{bacher2025joint}, normalised
  at the middle level, with $\eta_{up}$ estimated, not a separate age profile (Decision~\ref{od:eta}).\footnote{The PSID
  two-year log wage growth of continuously employed wives is 0.080, 0.046 and 0.027 at 25--34, 35--44 and 45--54
  (husbands 0.067, 0.029, 0.013; \code{IGE/60.PaperATargets.py}). At the start values, with $\pi^{up}_j$ rescaled to hold
  the growth at 25--34, $\eta_{up} = 0.5$ lowers the wage-growth block of the objective from 158 to 150, and
  $\eta_{up} = 1, 1.5, 2, 2.5$ to 145, 136, 123 and 109; at 0.5 the growth at 45--54 is still 0.045 (wives) and 0.037
  (husbands), at 2.5 the husbands' 0.032 against 0.013. The closed form from the mean levels by age would need
  $\eta_{up}$ of 6 (wives) to 16 (husbands), outside the bounds $[0, 3]$: the age profile of accumulation needs more than
  level dependence, and the remaining gap is a reported misfit. Start 0.5, with $\pi^{up}_j$ divided by the entry-level
  mean of $(h/4)^{-0.5}$ (1.051 wives, 1.019 husbands). Code: \code{IGE/paperA\_model/analytic\_start.py --ext hc\_eta},
  \code{model.py}; \code{results/diagnostics\_20261018.md}, step 4.}
  \item \emph{Displacement loss.} If $j$ is displaced at the end of the month, $h_j$ falls by
  $\lfloor d_j \rfloor$ levels and by one more with probability $d_j - \lfloor d_j \rfloor$, so that $d_j$ is the
  expected loss in levels.\footnote{The husband's labor income is 17\% (Black, $-\$4{,}160$, s.e.\ 757) and 24\% (White,
  $-\$9{,}622$, s.e.\ 1{,}071) lower in years 1--3 after his displacement and still about \$6{,}000--7{,}000 a year lower
  six years later, while his hours recover within four to five years: the lasting part is a wage loss. A wife's own
  displacement costs her about \$4{,}600 a year in labor income. Code: \code{IGE/19.FatherDisplacement.py},
  \code{IGE/25.JointJobLoss.py}. The pre-displacement dip in his hourly earnings ($-\$2{,}821$ in year $-1$, largest
  before plant closures) is reported, not modelled: there is no anticipation in the model, as in BPS, who find no
  advance information ($p=0.29$).} Whether $d_j$ is needed is decided by the earnings on re-employment by the length
  of non-employment after an involuntary separation (1--3, 4--12, 13--24 months; a target, Table~\ref{tab:targets}):
  depreciation alone makes the loss grow with the spell from close to zero, while $d_j$ adds a loss at re-employment
  after even a short spell. If the loss after 1--3 months is close to zero, $d_j$ is cut and the PSID years 1--3
  earnings loss must be reproduced by depreciation over the spell; if it is not, $d_j$ stays.\footnote{For all
  separations in the SIPP 1995--2016, \citet{bacher2025joint} (their Table 8) report earnings in the first month of
  re-employment 1\%, 7\% and 20\% below the last month of the previous job after 1--3, 4--12 and 13--24 months of
  non-employment, and match it with depreciation that does not depend on the level. Paper A's version restricts the
  separations to involuntary ones and reports it by spouse, from the SIPP panel 2003--2016 (\code{IGE/40.SIPPJointLoss.py};
  computed in \code{IGE/60.PaperATargets.py}).}
  \item \emph{Persistent innovations.} With probability $p^c$ a shared innovation $\zeta_c \in \{-1, +1\}$ (equally
  likely) moves both spouses' levels in the same direction; independently, with probability $p^\varepsilon_j$ an
  idiosyncratic innovation $\zeta_j \in \{-1, +1\}$ moves spouse $j$'s level. After a separation the wife's level
  receives innovations with probability $p^c + p^\varepsilon_w$.
\end{enumerate}
The innovations are the permanent wage shocks of BPS's wage process, represented on the human capital grid: with
annual horizons and no truncation, the annual variance of the permanent component of spouse $j$ is
$12\,\psi_j^2\,(p^c + p^\varepsilon_j)$ and the covariance across spouses $12\,\psi_w\psi_h\,p^c$. A seven-point grid is
not a unit root; over the ages BPS use (30--57) it is close to one if the grid is wide, which is why $\psi_j$ is set
from the cross-sectional dispersion (Table~\ref{tab:params}).\footnote{The decision to give the spouses a shared wage
shock with a persistent and a transitory component was taken for the lifecycle model on PSID evidence (her wage change
on his earnings change, Black 0.091$^{***}$, White 0.048$^{***}$; \code{ReducedForm/51}, \code{52}); BPS's Table 3, which
the reconciliation reproduces in all 36 cells, gives a transitory covariance of 0.0058$^{**}$ (correlation 0.31) and a
permanent covariance of 0.0027 (correlation 0.08, $p=0.246$). Joint displacement years raise the covariance of the
spouses' \emph{earnings} residuals ($+0.193^{**}$, $p=0.019$), not of their wage residuals (0.0063$^{**}$ without any
displacement year, 0.0054$^{**}$ without joint ones): the joint loss is an employment shock, and it is modelled as one,
separately from the shared wage shocks. Code: \code{IGE/bps\_reconciliation/} (\code{bps\_reconciliation.md}).}

\paragraph{Transitory shocks: one draw per calendar year.} The transitory shocks $(\nu_w, \nu_h)$ are drawn in January,
at the start of the month and before the hours choice, and held for the twelve months of the calendar year; they are
independent across years. Each spouse's shock takes two values, $\pm\sigma_{\nu,j}$, and the pair has probabilities
$(1+\rho_\nu)/4$ for equal signs and $(1-\rho_\nu)/4$ for opposite signs, so that the variances are $\sigma_{\nu,j}^2$
and the correlation $\rho_\nu$. A separated wife keeps her own component. The shock is therefore a state within the year
(Table~\ref{tab:state}). The reasons for an annual rather than a monthly draw (Decision~\ref{od:trans}) are what each
source measures.
\begin{enumerate}
  \item In \citet{blundell2016consumption} the transitory shock is the serially uncorrelated component of the
  \emph{annual} hourly wage (annual earnings divided by annual hours), estimated from two-year differences of biennial
  data and described as short-lived events such as a short illness; it is assumed fully smoothed. Its variance is 0.0275
  for husbands and 0.0125 for wives, with a correlation of 0.31 across spouses. Part of it is measurement error, which
  hourly wages inherit from both earnings and hours (division bias); BPS net out the error with the shares of the
  validation study of \citet{bound1994evidence}.
  \item \citet{bacher2025joint} have no transitory wage shock: the wage is a share (0.7) of the human capital level,
  which moves one grid step up with some probability in each month of work and one step down with some probability in
  each month out of work. Monthly earnings risk comes from job loss and human capital.
  \item \citet{ganong2026earnings} show that monthly earnings within a job are volatile while hourly wages are largely
  stable from month to month: the volatility is in hours, and largely set by firms.
  \item A monthly i.i.d.\ shock that reproduces BPS's annual variance through the annual average of monthly wages needs
  a monthly standard deviation of about $\sqrt{12 \times 0.0275} = 0.57$ for husbands ($\sqrt{12 \times 0.0125} = 0.39$
  for wives): a typical monthly swing of 57\% in the hourly wage, which contradicts point 3. An annual draw carries BPS's
  object with standard deviations of at most $\sqrt{0.0275} = 0.166$ and $\sqrt{0.0125} = 0.112$, the same in every month
  of the year.
\end{enumerate}
Month-to-month risk in the model therefore comes from displacement and job finding, not from the wage rate. The model's
true transitory variances are smaller than BPS's estimates, because part of BPS's transitory variance is measurement error
that the model does not contain: that part is added only in the measurement layer of the BPS test, as the two-component
error mixture of Section~\ref{sec:bps}, and is not a shock the household faces. In practice $(\sigma_{\nu,w},
\sigma_{\nu,h}, \rho_\nu)$ start at BPS's values (an upper bound) and are reduced, once and before estimation, by the
amount that BPS's estimator, run on the simulated BPS-style panel with the error mixture at the start values, attributes
to transitory variance from the model's other sources (displacement losses and human capital moves within a year and
the mixture); they are then held fixed. The cost of the annual draw is a four-point state; the work per month is of the
same order as with nine quadrature nodes for a monthly shock (Section~\ref{sec:solution}).

\subsection{Displacement risk in three parts}

A job ends involuntarily only through a displacement, which in the model is a permanent job loss.\footnote{Temporary
layoffs with recall are outside the model. In the CPS they account for about a third of the same-month excess: without
the husband's temporary layoffs the excess is 2.7 points (0.027$^{***}$), and without both spouses' 2.6 points on a base
of 0.23\%; the strictest joint permanent loss (both still unemployed two months later, outside temporary layoffs, the
summer and the pandemic) is 13 times as likely in his loss month (1.6 points on 0.17\%). The CPS targets are therefore
computed on non-temporary job losses (Decision~\ref{od:recall}). \citet{bacher2025joint} find their added-worker effect
of similar size for permanent separations only (their Online Appendix Table B.4). Code:
\code{IGE/37.CPSJointLossChecks.py}.} At the end of month $t$ the displacement shocks are drawn in two stages:
\begin{enumerate}
  \item \emph{Common shock.} With probability $\pi_c(s)$, which depends on the workplace match ($s = 0$: different
  industries; $s = 1$: same detailed industry; $s = 2$: the shared-employer cell), a common shock hits the couple and
  every employed spouse is displaced.
  \item \emph{Own hazards.} Otherwise each employed spouse is displaced independently with probability
  $\delta_j(r)\, m_\tau$, where $m_L = 1$ and $m_H > 1$ is the risk multiplier of high-risk couples, whose share
  is $p_H$.
\end{enumerate}
When both spouses work, the probability that both are displaced in the same month is
\[
  \Pr(\text{both}) = \pi_c(s) + \big(1 - \pi_c(s)\big)\, m_\tau^2\, \delta_w(r)\,\delta_h(r),
\]
and the probability that only spouse $j$ is, $\big(1 - \pi_c(s)\big)\, m_\tau\delta_j(r)\,\big(1 - m_\tau
\delta_{-j}(r)\big)$. The three parts are separately identified because they leave different traces at the monthly
frequency: the common shock produces excess joint loss only in the month it hits, the couple type a flat elevation of
the wife's hazard in every month around the husband's loss and in the PSID placebo window, and the own hazards the base
rates (Section~\ref{sec:ident}). The monthly elevation pins the couple type along a ridge rather than at a point: a rare
and very risky type and a common and mildly risky one give the same CPS leads and lags, but the rare type gives the
smaller PSID same-year excess. The estimation starts on the ridge at $p_H = 0.02$, $m_H = 15.95$, and estimates both
(Decision~\ref{od:ridge}).\footnote{Along the ridge (own hazards rescaled to hold the base rates) the PSID excess falls
from 0.118 at $p_H = 0.10$ to 0.071 at $p_H = 0.02$ (data 0.063, s.e.\ 0.005), while the CPS leads and lags stay at
0.0034--0.0061 (data 0.0056--0.0067); off the ridge a smaller $m_H$ at $p_H = 0.10$ fits the annual excess only by
losing the leads and lags ($m_H = 4$: 0.0015). Joint-loss and base-hazard blocks of the objective: 401 at $p_H = 0.10$,
185 at 0.02. Tested and rejected: a short memory of the common shock, each spouse's hazard times $1 + c_{mem}$ while the
other is non-employed and entitled after a displacement (one parameter, $c_{mem} \in [0, 50]$): at its start
($c_{mem} = 5.82$, $m_H = 4.33$) the joint-loss blocks are 177 against 62 on the ridge, and the excess one month before
and after his loss 0.013--0.014 against 0.006. It stays as an option (\code{estimate.py --cs-memory on}, off by
default). Code: \code{IGE/paperA\_model/results/diagnostics\_20261018/profile\_sigma\_ridge.py},
\code{analytic\_start.py} (\code{P\_H\_START}); \code{results/diagnostics\_20261018.md}, B and steps 2a, 2.} The model
has no aggregate state; the larger excess in recessions is used in a comparative static (Section~\ref{sec:cf}).\footnote{The CPS excess is 6.7 points in recessions against 3.6 points
otherwise (difference 0.031$^{***}$); the change in the state unemployment rate does not moderate the PSID excess
(0.020, s.e.\ 0.017). Code: \code{IGE/30.CPSJointJobLoss.py}, \code{IGE/25.JointJobLoss.py}.}

\subsection{The search friction}

A spouse who is non-employed at the start of month $t$ finds a job by the end of the month with probability
$\lambda_j(r)$; this probability measures the search friction. It depends on race and spouse only, not on age, the human
capital level or the length of the spell. The known risk of that restriction is age: in \citet{bacher2025joint}, where
job finding is set by firms, the average rate for non-employed spouses of one-earner couples falls from about 0.16 a
month at 26 to about 0.01 at 64 (their Figure 10), mostly through the age distribution of assets and human capital, so a
constant $\lambda_w$ may overstate the re-entry of older wives. The guard is a validation by age: the CPS monthly entry
rate of non-employed wives and the SIPP re-employment of displaced husbands by age group, and job finding by age against
Bacher et al.'s; one age slope in $\lambda_j$ is added only if that validation fails (Decision~\ref{od:husband}). A
spouse displaced at the end of month $t$ is non-employed in $t+1$ and can find a job from the end of $t+1$, the
convention of \citet{bacher2025joint} (their footnote 10). A job allows either hours level. A wife without UI
entitlement who finds a job chooses at the start of the next month whether and how much to work, and declining it
returns her to non-employment. A wife who is entitled decides at the end of the month, after both spouses'
displacement and job-finding events and before the human capital, child and separation draws, whether to accept, as
in \citet{bacher2025joint}: accepting ends her entitlement; declining keeps it, and it then lapses only at the monthly
hazard $1/6$. The decision is binary, with mean-zero extreme-value shocks of scale $\sigma_\varepsilon$, and adds no
state (\code{solve.py}, the offer node). A wife can also leave a job she holds; her monthly exit from employment
outside a job loss is a target (Table~\ref{tab:targets}), so that margin is disciplined by data rather than left
free.\footnote{\citet{bacher2025joint} (their Table 6) report pooled monthly transitions of 1.8\% from employment to
non-participation and 0.9\% from employment to unemployment, for all reasons. Paper A's target is wives with an employed
husband, without a job-loss reason, from the CPS panel 1994--2023 (\code{IGE/30.CPSJointJobLoss.py}; computed in
\code{IGE/60.PaperATargets.py}).} The husband always takes a job he finds.\footnote{Among wives with no job
the month before, looking rises at once and employment builds up over about six months after the husband's separation
(SIPP), and in the CPS wives who were already unemployed find work less often after his loss ($-0.020^{*}$ after one
month, $-0.050^{**}$ after three). The second fact suggests that the common shock also lowers her job finding; it is
kept as validation rather than modelled (parsimony). Code: \code{IGE/40.SIPPJointLoss.py},
\code{IGE/30.CPSJointJobLoss.py}.}

% =====================================================================================================
\section{The government and the budget constraint}\label{sec:gov}

The government is represented by its schedules, because the paper's policy question is what each schedule pays in the
state in which no one earns.\footnote{BPS's tax-and-transfer function is log-linear in positive family earnings and
contains the federal income tax, the EITC and food stamps, but no unemployment insurance (their
\code{tax\_rate\_estimation\_snap.do}); the state in which no one earns is outside its domain. Code reading in
\code{IGE/proposals/fable\_bps2016\_paperA.md}, Section 2.2.} All schedules are set by policy era $e$, of which there
are three (Decision~\ref{od:eras}).

\paragraph{Income and payroll taxes.} On positive annualised income $Y^a = 12\,(y_w + y_h + b_w + b_h)$ the household
keeps $(1 - \chi_e)\,(Y^a)^{1 - \mu_e}$, BPS's functional form, with $(\chi_e, \mu_e)$ re-estimated by era from the
federal income and payroll taxes of married couples filing jointly, without the EITC and SNAP, which enter separately.
The monthly tax is $T_e = \big[Y^a - (1-\chi_e)(Y^a)^{1-\mu_e}\big]/12$ (zero when $Y^a = 0$).

\paragraph{Unemployment insurance.} A displaced spouse becomes entitled to UI with probability $q_{UI}$ (take-up and
monetary eligibility together); a spouse who leaves a job is not entitled. Each spouse has one UI state, $u_j \in
\{0, 1\}$ (not entitled, entitled). An entitled spouse who is still non-employed loses entitlement at the end of each
month with probability $1/6$, so that expected duration is six months, the 26 weeks of regular benefits; this
geometric expiry is that of \citet{bacher2025joint} and replaces a duration clock.\footnote{The geometric hazard ends
a sixth of entitlements in the first month and leaves a third beyond six months ($(5/6)^6 = 0.33$), where the statute
ends all at 26 weeks. The moments that see UI duration---the SIPP re-employment paths at 3, 6 and 12 months and the
PSID annual moments---cannot tell a geometric expiry from a two-stage clock with the same mean (three months, then
three more, each stage ending with probability $1/3$), and the geometric form needs half the labor states. Counterfactual (ii)(b) sets the hazard to $1/12$ (52 weeks).} While $u_j = 1$ the benefit is
\[
  b_j = \min\big\{\rho_{UI}\; \bar w_j(h_j)\, H(\ell_j^{lost}),\; \bar b_{e}\big\},
  \qquad \bar w_j(h) = \exp\{\mu_j(r) + \psi_j(h - 4)\},
\]
with replacement rate $\rho_{UI} = 0.5$, cap $\bar b_e$ the mean of the states' maximum weekly benefits in the era
(monthly), and $\ell_j^{lost}$ the hours of the lost job, which the wife's UI state records ($P$ or $F$) and which is
$F$ for the husband. The benefit is computed at the current level, after the displacement loss; the error is at most
one step and is the same in every counterfactual arm. The wife's labor state therefore takes four values (employed;
non-employed and entitled after a part-time job; entitled after a full-time job; not entitled) and the husband's three
(employed; entitled; not entitled). The lost-hours record is on the cut list (Decision~\ref{od:count}): the alternative is
a benefit on full-time-equivalent earnings times a part-time factor.

\paragraph{Earned Income Tax Credit.} The credit is paid monthly on annualised family earnings:
\begin{gather*}
  \mathcal E = \tfrac{1}{12}\,\mathrm{EITC}_e\big(12\,(y_w + y_h);\, n\big), \\
  \mathrm{EITC}_e(Y; n) = \max\Big\{0,\; \min\{\epsilon^{in}_{e,n} Y,\, \bar E_{e,n}\} - \epsilon^{out}_{e,n}\max\{0,\, Y - Y^{out}_{e,n}\}\Big\},
\end{gather*}
with phase-in rate, maximum credit, phase-out rate and threshold for married filers with $n$ children. Paying on
annualised current earnings rather than carrying an ``earnings so far this year'' state is the standard approximation
in monthly models; the counterfactual of Section~\ref{sec:cf} compares schedules, so it applies to both arms. The
credit is zero when no one earns.

\paragraph{SNAP.} A household of size $n_{hh}$ ($2 + n$ married, $1 + n$ separated) is eligible when gross income
$Y^g = y_w + y_h + b_w + b_h$ is at most 130\% of the poverty line and assets are below the asset limit of the era. Its
benefit is
\[
  B_S = \max\big\{0,\; A_e(n_{hh}) - 0.3\,\max\{0,\; Y^g - 0.2\,(y_w + y_h) - D_e(n_{hh}) - c_{cc}\}\big\},
\]
the maximum allotment minus 30\% of net income after the 20\% earned-income deduction, the standard deduction and the
dependent-care deduction $c_{cc}$. Take-up $S \in \{0, 1\}$ is chosen (with the utility cost $\kappa_S$).

\paragraph{Childcare.} The wife's work costs $c_{cc} = p_{cc}(k)\, H(\ell_w)$, with an hourly price that depends on the
youngest child's age bucket. A subsidy lowers $p_{cc}$.

\paragraph{Number of children.} The EITC and SNAP depend on the number of children, which is not a state: it is set to
$n = n(k, \text{age}_w, r)$, the PSID mean for couples with the same youngest-child bucket, wife's age group and race,
rounded to one or two.

\paragraph{Consumption floor.} A transfer $Tr \ge 0$ guarantees spending of at least $c_{min}$ (equivalised), standing
for the transfers that are not modelled (AFDC/TANF, which barely moves after displacement).

\paragraph{Budget constraint.} Cash on hand of a married couple is
\begin{equation}\label{eq:budget}
  \begin{aligned}
  X^M = {} & \big(1 + r_m + \omega\,\1\{a < 0\}\big)\,a + y_w + y_h + b_w + b_h - T_e + \mathcal E \\
           & + S\, B_S - c_{cc} - F\,\big(\1\{\ell_w > 0\} + \1\{\ell_h > 0\}\big) + Tr,
  \end{aligned}
\end{equation}
and the household chooses $c$ and next month's assets $a' = X^M - c \ge -\underline{a}$. The monthly interest rate is
$r_m = 1.03^{1/12} - 1$ on saving; on debt it is $r_m + \omega$, where the wedge
$\omega = 1.114^{1/12} - 1.03^{1/12} = 0.0066$ a month makes the borrowing rate the real interest rate on credit cards,
11.4\% a year, 8.4 points above the saving rate (Decision~\ref{od:wedge}). The wedge is a coded input, not a parameter;
the asset grid carries a node at zero, where the budget kinks.\footnote{Federal Reserve G.19, interest rate on credit
card plans of commercial banks, all accounts (FRED \code{TERMCBCCALLNS}): 13.7\% nominal on average over the 108
quarters of 1995--2021, less 12-month CPI-U inflation (2.3\%), less the real saving rate of 3\%; accounts assessed
interest give 12.2\% real, the 1997--2008 quarters 11.2\%. The targets' liquid wealth nets unsecured debt against liquid
assets, and negative liquid wealth is mostly debt (\code{ReducedForm/62.PoorMarriedCouples.py}), so the credit-card rate
is the marginal rate on it. At the start values and on common draws the wedge lowers the objective from 68{,}086 to
64{,}657 and the share of wives aged 25--34 in couples with liquid wealth at or below zero from 0.80 to 0.48 (data
0.57). It lowers the objective at every borrowing limit tried ($\underline a$ = 15, 25.6, 40 thousand dollars: 65{,}618,
64{,}657, 63{,}529 with the wedge against 66{,}522, 68{,}086, 69{,}845 without), and with the wedge the wealth block
falls from 3{,}079 to 2{,}672, 2{,}419 and 2{,}317 as $\underline a$ rises from 15 to 25.6, 40 and 55. The shape at
45--59 stays wrong with or without it: too few couples at or below zero (0.085 against 0.381) and too many above a year
of earnings (0.388 against 0.177), an open misfit. Code: \code{IGE/paperA\_model/borrowing\_wedge.py},
\code{inputs/borrowing\_wedge.json}; \code{results/diagnostics\_20261018.md}, step 3.} The separated wife's cash on hand $X^S$ is the same expression
without the husband's terms, with the single-filer EITC and SNAP schedules and the childcare price of her child.

% =====================================================================================================
\section{Marriage}\label{sec:marriage}

The marriage ends at the end of a month with probability
\begin{equation}\label{eq:sep}
  \sigma(r, e_h') = \sigma_0(r) + \sigma_J(r)\, \1\{e_h' = 0\},
\end{equation}
which is higher while the husband is out of work after a displacement ($e_h' = 0$ is his status at the start of next
month, and he is non-employed only after a displacement). At a separation assets are divided equally; the wife continues
alone with the children, with her own labor state, human capital, UI state, transitory shock and child state, and the
husband leaves the model (Decision~\ref{od:welfare}). There is no remarriage within the horizon and no child support. The separation is exogenous (parsimony): the
model reproduces the separation response of Paper A, Section 5, and its difference by race; the household anticipates
the hazard, which adds a precautionary motive.\footnote{The share of displaced fathers still living with the wife they
had before the job loss falls by 4.5 points for Black fathers ($-0.045^{***}$, s.e.\ 0.015) and 1.9 points for White
fathers ($-0.019^{**}$, s.e.\ 0.009) over years 1--3, and by 7.4 (s.e.\ 2.1) and 3.0 points (s.e.\ 1.4) by year 6. In
BPS's panel a separation follows a wife's displacement 2.9 points more often ($p=0.033$) and the husband's 1.0 point
more often ($p=0.26$). Code: \code{IGE/19.FatherDisplacement.py}, \code{IGE/bps\_reconciliation/}. Tying the hazard to his
non-employment adds no state variable (Decision~\ref{od:sep}); a twelve-month flag after his displacement (state
$\times 2$) and a hazard that also responds to her displacement were considered and not adopted.} The hazard is a
reduced form for the marital surplus. The marriage is kept exogenous because the data pin the mean response and little
that would discipline a decision: the response is the same after plant closures as after layoffs, which is what a fall
in household income and surplus predicts rather than news about the husband \citep{charles2004job}, and it does not
differ significantly by race, by the child's age or by the wife's employment.\footnote{Still with the baseline wife,
years 1--3: closures minus layoffs 0.000 ($p = 0.99$, Black) and 0.006 ($p = 0.75$, White); by the youngest child's age
$-0.030$ and $-0.030$; wife employed at baseline $-0.036^{***}$ against not employed $-0.019$ (difference 0.017, s.e.\
about 0.017); the race difference in the targeted mean is 0.5 standard errors. The outcome is one by construction
before the displacement, so it has no placebo test. The gap keeps widening after he is re-employed: $-0.027$, $-0.037$
and $-0.050$ in years 1, 3 and 6 (Black $-0.037$, $-0.053$, $-0.074$; White $-0.017$, $-0.022$, $-0.030$). Code:
\code{IGE/19.FatherDisplacement.py}, \code{IGE/38.PSIDJointLossChecks.py}; \code{IGE/proposals/fable\_marriage\_divorce\_paperA.md}.}
Its time path is an untargeted test of the form: the hazard is raised only while he is out of work, whereas the PSID
gap keeps widening after he is back at work. The husband's welfare after a separation is not counted.

% =====================================================================================================
\section{Timing within the month}\label{sec:timing}

The simulator runs each month in this order; the Bellman equations of Section~\ref{sec:problem} follow the same order.
\begin{enumerate}
  \item \textbf{State.} The month begins with the state $\Omega_t$ (Section~\ref{sec:problem}) known.
  \item \textbf{Transitory wage shocks}: in January, $(\nu_w, \nu_h)$ are drawn and observed (only $\nu_w$ for a
  separated wife); in the other months they keep their January values.
  \item \textbf{Taste shocks} $\varepsilon(d)$ are drawn for every feasible alternative $d = (\ell_w, \ell_h, S)$:
  $\ell_w \in \{0, P, F\}$ if the wife has a job, $\ell_w = 0$ otherwise; $\ell_h \in \{P, F\}$ if the husband has a
  job; $S \in \{0, 1\}$ if the household would be SNAP-eligible under $(\ell_w, \ell_h)$.
  \item \textbf{Choice.} The household chooses $d$ and then consumption and saving given $d$.
  \item \textbf{Income.} Earnings, UI, taxes, EITC, SNAP, childcare and work costs are paid; consumption takes place.
  \item \textbf{Labor-market shocks} at the end of the month, in this order: (a) the common shock and the own
  displacement hazards for spouses who worked this month; (b) job finding for spouses who were non-employed this month,
and an entitled wife's decision to accept or decline the job (Section~\ref{sec:labor});
  (c) a wife who chose $\ell_w = 0$ while holding a job becomes non-employed without UI entitlement.
  \item \textbf{Human capital} moves in the three steps of Section~\ref{sec:labor} (accumulation or depreciation
  according to this month's work, the displacement loss, the persistent innovations).
  \item \textbf{UI states}: a spouse displaced this month becomes entitled with probability $q_{UI}$ (recording the
  wife's hours); an entitled spouse who was non-employed this month and did not find a job loses entitlement with
  probability $1/6$; a spouse who found a job leaves UI.
  \item \textbf{Children}: the youngest-child bucket moves (Section~\ref{sec:problem}).
  \item \textbf{Marriage}: the separation is drawn with probability $\sigma(r, e_h')$; at a separation assets are
  halved and the wife continues in the separated problem.
\end{enumerate}
A displacement at the end of month $t$ is a job loss in month $t$ in the data's sense (employed in $t$, not employed in
$t+1$), so a common shock produces both spouses' job losses in the same CPS month. A displaced spouse cannot find a job
in the month of the displacement, as in \citet{bacher2025joint} (their footnote 10).

% =====================================================================================================
\section{The household's problem}\label{sec:problem}

\paragraph{State.} The married couple's state at the start of month $t$ is
\[
  \Omega_t = \big(t,\; a,\; e_w, u_w,\; h_w,\; e_h, u_h,\; h_h,\; k,\; \nu;\; s, \tau, r\big),
\]
with assets $a$, employment $e_j \in \{0, 1\}$ (holding a job), the UI states ($u_j = 0$ when employed; the wife's
records the hours of the lost job), the two human capital levels, the youngest-child bucket $k$, this year's transitory
shocks $\nu = (\nu_w, \nu_h)$, and the permanent characteristics. The separated wife's state is
$\Omega^S_t = (t, a, e_w, u_w, h_w, k, \nu_w; \tau, r)$. Table~\ref{tab:state} lists the grids. The permanent
characteristics $(r, s, \tau)$ are known and never change, so each combination is a separate problem: it is solved and
its couples simulated before the next one is solved.

\begin{table}[htbp]\centering\small
\caption{State space (married couple, one month)}\label{tab:state}
\begin{tabular}{L{4.2cm}L{5.2cm}cL{3.6cm}}
\toprule
State & Values & Points & Why \\
\midrule
Assets $a$ & curved grid, dense near the borrowing limit & 30 & savings, liquidity \\
Wife's labor state $(e_w, u_w)$ & employed; non-employed entitled to UI after a part-time job; after a full-time job; not entitled & 4 & UI per spouse \\
Husband's labor state $(e_h, u_h)$ & employed; non-employed entitled; not entitled & 3 & UI per spouse \\
Wife's human capital $h_w$ & 1--7 & 7 & wage, displacement loss \\
Husband's human capital $h_h$ & 1--7 & 7 & wage, displacement loss \\
Youngest child $k$ & none, 0--5, 6--17 & 3 & work cost by child age \\
Transitory shocks $\nu$ & $\pm\sigma_{\nu,w} \times \pm\sigma_{\nu,h}$, drawn in January & 4 & BPS's annual shock \\
\midrule
Total per month, one problem & $30 \times 4 \times 3 \times 7 \times 7 \times 3 \times 4$ & 211{,}680 & \\
\midrule
Workplace match $s$ & different industry; same detailed industry; shared-employer cell & 3 & common shock \\
Couple type $\tau$ & low, high risk & 2 & assortative risk \\
Race $r$ & Black, White & 2 & environment \\
\midrule
Total per month, all problems & $211{,}680 \times 3 \times 2 \times 2$ & 2{,}540{,}160 & \\
\bottomrule
\end{tabular}
\begin{minipage}{0.95\textwidth}\footnotesize The age $t$ (420 months) and the policy era index the solution. The
separated wife's state has $30 \times 4 \times 7 \times 3 \times 2 = 5{,}040$ points a month for each $(\tau, r)$.
Without the transitory state (a monthly i.i.d.\ shock, not adopted) a problem would have 52{,}920 points.\end{minipage}
\end{table}

\paragraph{Married couple.} Write this year's shocks $\nu = (\nu_w, \nu_h)$ separately from the other states $\Omega$
and let $\mathcal D(\Omega, \nu)$ be the feasible alternatives. The value at the start of the month, after $\nu$ is
observed and before the taste shocks, is
\begin{equation}\label{eq:VM}
  V^M_t(\Omega, \nu) = \E_\varepsilon \max_{d \in \mathcal D(\Omega,\nu)} \big\{ v^M_t(d; \Omega, \nu) + \sigma_\varepsilon\,\varepsilon(d) \big\}
  = \sigma_\varepsilon \ln \sum_{d \in \mathcal D(\Omega,\nu)} \exp\!\big(v^M_t(d; \Omega, \nu)/\sigma_\varepsilon\big) + \sigma_\varepsilon\gamma_E ,
\end{equation}
with $\gamma_E$ Euler's constant. The choice-specific value is
\begin{equation}\label{eq:vM}
  v^M_t(d; \Omega, \nu) = \max_{a' \ge -\underline a} \Big\{ u\!\Big(\frac{X^M(d; \Omega, \nu) - a'}{s_M\, e(k)}\Big)
     - \phi_w(\ell_w, k) - \phi_h(\ell_h) - \kappa_S S + \beta\, W^M_{t+1}(a'; d, \Omega, \nu) \Big\},
\end{equation}
and the continuation value integrates over next month's state and shocks:
\begin{equation}\label{eq:WM}
  W^M_{t+1}(a'; d, \Omega, \nu) = \sum_{\Omega'} P(\Omega' \mid \Omega, d, a')
   \Big[ \big(1 - \sigma(r, e_h')\big)\, \E_{\nu' \mid \nu} V^M_{t+1}(\Omega', \nu')
         + \sigma(r, e_h')\, \E_{\nu_w' \mid \nu_w} V^S_{t+1}\big(\Omega'^S, \nu_w'\big) \Big],
\end{equation}
where $\Omega'^S$ keeps the wife's components of $\Omega'$ with assets $a'/2$. The expectation over the shocks is degenerate within the year
($\nu' = \nu$) and is the four-point sum over a new draw when $t+1$ is a January. The transition $P$ factors into the
events of Section~\ref{sec:timing}: displacement (common shock, then own hazards, both depending on $s$, $\tau$, $r$ and
on who worked, $\ell_j > 0$), job finding $\lambda_j(r)$, the wife's quit when $\ell_w = 0$, the three human capital
steps (whose probabilities depend on $\ell_j$ and on the displacement), the UI states ($q_{UI}$, the expiry hazard $1/6$
and the wife's lost hours), and the child bucket. Conditional on the month's displacement and job-finding events, the
human capital, UI and child transitions are independent, so the expectation is taken one dimension at a time.

\paragraph{Separated wife.}
\begin{align}
  V^S_t(\Omega^S, \nu_w) &= \sigma_\varepsilon \ln \sum_{d \in \mathcal D^S} \exp\!\big(v^S_t(d; \Omega^S, \nu_w)/\sigma_\varepsilon\big) + \sigma_\varepsilon\gamma_E, \label{eq:VS}\\
  v^S_t(d; \Omega^S, \nu_w) &= \max_{a' \ge -\underline a} \Big\{ u\!\Big(\frac{X^S(d; \Omega^S, \nu_w) - a'}{e(k)}\Big)
     - \phi_w(\ell_w, k) - \kappa_S S \nonumber \\
     &\qquad\qquad + \beta \sum_{\Omega^{S\prime}} P^S(\Omega^{S\prime} \mid \Omega^S, d, a')\, \E_{\nu_w' \mid \nu_w} V^S_{t+1}(\Omega^{S\prime}, \nu_w') \Big\}, \nonumber
\end{align}
with $d = (\ell_w, S)$, her own displacement hazard $\delta_w(r)\,m_\tau$ (the common shock no longer applies) and job
finding $\lambda_w(r)$.

\paragraph{Children.} The bucket $k$ follows a monthly Markov chain estimated outside the model on PSID couples: from
$k=0$ a birth moves it to $k=1$ with probability $b_0(\text{age}_w, r)$; from $k=1$ it moves to $k=2$ with probability
$1/72$ unless a birth occurs; from $k=2$ a birth returns it to $k=1$ with probability $b_2(\text{age}_w, r)$ and it
moves to $k=0$ with probability $1/144$. The same stochastic ageing is used in the solution and in the simulation, so
that the household's expectations are correct for the process it faces; a finer child-age grid is a refinement if the
added-work moments by child age turn out to be sensitive.

\paragraph{Terminal value.} At the wife's age 60 the couple retires: $V^M_{T+1}(a, h_w, h_h)$ is the discounted utility
of 240 months of spending equal to the annuity value of $a$ plus Social Security benefits of 40\% of each spouse's
earnings at its final level, divided by $s_M$; $V^S_{T+1}(a, h_w)$ is the same for the wife alone.

\paragraph{Initial conditions.}\label{sec:init} Each simulated couple is a copy of a PSID couple, entering in January
of the year of its first observation from 1984 on (couples observed earlier enter with their 1984 state) with its
calendar year, the wife's age, race, workplace match, liquid wealth (imputed by age, race and education in waves without wealth), each spouse's
employment, the youngest child's bucket, and each spouse's human capital level, set from the log hourly wage at entry
(for a non-employed spouse, the last observed wage or the education-and-age prediction). The couple type is drawn
with probability $p_H$, independently of the observables, and the first year's transitory shocks from their
distribution. A UI state at entry is set to entitled for a spouse who reports a job loss in the entry year and is not
employed, and to not entitled otherwise. Each PSID couple is copied $K$ times with independent draws.

% =====================================================================================================
\section{The children block}\label{sec:children}

The children block is outside the estimation; the parents' model is estimated without it, so its results for
consumption and the ranking of policies stand on their own. For each child present in a simulated household it maps
family income by childhood age and the presence of both parents into an adult outcome,
\begin{equation}\label{eq:child}
  \ln Y^{adult}_i = \alpha + \beta_{05} \ln \bar y_{i,0\text{--}5} + \beta_{617} \ln \bar y_{i,6\text{--}17}
    + \beta_P\, P_i + \eta_i,
\end{equation}
where $\bar y_{i,g}$ is mean family income (earnings, UI, EITC and SNAP) over the months in which the youngest child is
in age group $g$, and $P_i$ the share of those months with both parents present. A counterfactual changes adult
outcomes by $\beta_{05}\Delta\ln\bar y_{0\text{--}5} + \beta_{617}\Delta\ln\bar y_{6\text{--}17} + \beta_P\Delta P$. The slopes
are not estimated. Two scenarios bracket them: a high scenario in which a father's displacement during childhood, as
simulated, lowers adult earnings by about 9\%, the estimate in Canadian administrative data
\citep{oreopoulos2008intergenerational} and the PSID association of \citet{page2009parental}; and a low scenario at the
smaller effects in US tax data \citep{hilger2016parental}. The PSID associations of Paper A, Section 6, are reported next
to them as an upper bound only.\footnote{A displacement before age six is associated with a household income rank 3.9
percentiles lower and 0.28 fewer years of schooling; one at 12--17 with adult earnings about 14\% lower. These are
associations that bound, rather than identify, the effect: displaced families' income was 24 log points lower before
the job loss, and with plant closures alone the associations vanish for Black children. Code:
\code{IGE/20.PaperDisplacement.py}, \code{IGE/22.DisplacementID.py}. Siblings are not used as controls.}
The model's task is to translate the parents' outcomes into the children's, not to estimate the technology of child
development (Decision~\ref{od:child}). The parental-presence term is kept because the separation of the couple is one
of the channels by which the joint loss reaches the children (the separation response of Section~\ref{sec:marriage}):
without it, a counterfactual that changes separations would change nothing for the children except through income. Its
slope $\beta_P$ is set in the high scenario at the PSID association of the share of childhood months with both parents
present with adult earnings, computed with the controls of Paper A, Section 6 (\code{IGE/20.PaperDisplacement.py}, to be
added), which bounds it from above, and at zero in the low scenario.

% =====================================================================================================
\section{Solution and simulation}\label{sec:solution}

\paragraph{Solution.} Backward induction over 420 months for each problem $(r, s, \tau)$ and policy era. Within a month
the discrete alternatives are handled with extreme-value taste shocks and consumption with the endogenous grid method
for discrete-continuous models (DC-EGM; \citealp{carroll2006method, iskhakov2017endogenous}), with the upper envelope
over the alternatives, as in \citet{bacher2025joint}. Between asset nodes the continuation value is a monotone cubic
Hermite interpolant of the node values and their exact derivatives (the derivative the Euler step inverts), and the
endogenous-grid point is polished by Newton steps on the Hermite first-order condition where consumption changes by
more than 5\% between bracketing points.

\paragraph{The SNAP asset test.} The asset test puts a notch in next month's value at the asset limit $L$ (1{,}884
dollars in the 1997--2008 schedule): a couple whose assets cross it loses the SNAP alternative. The asset grid carries
nodes at $L$ and $L + 10^{-5}$ thousand dollars, so that the drop lies inside one segment and the continuation value is
smooth on every other; the corner $a' = L$ is a candidate in every state; where the endogenous grid folds at the drop,
a two-branch endogenous grid is used, with a nine-point search and up to six Newton steps where the marginal value
rises.\footnote{Without this treatment the Hermite continuation value smeared the drop and the first-order conditions
missed the bunching at $a' = L$: the gap between the endogenous-grid solution and a brute-force maximisation over $a'$
with the same continuation values reached 0.16 (99th percentile $3.4 \times 10^{-3}$) at low cash on hand. With it
the 99th percentile and maximum are $1.7 \times 10^{-4}$ and $2.6 \times 10^{-3}$ at the start values,
$7.0 \times 10^{-4}$ and $2.4 \times 10^{-3}$ with $\kappa_S = 0$, and $3.7 \times 10^{-5}$ and $1.2 \times 10^{-4}$
with $\kappa_S = 2.64$ (criteria $10^{-3}$ and $5 \times 10^{-3}$); 79\% of states keep a monotone grid or one monotone
on both sides of $L$. Code: \code{model.py} (\code{ASSET\_EPS}, \code{jL}), \code{solve.py};
\code{results/diagnostics\_20261018.md}, follow-up 1.}

\paragraph{Estimation size.} The model is estimated at a reduced size: seven human capital levels, the workplace groups
$s = 1$ and $s = 2$ merged (both use $\pi_c(1)$), the 1997--2008 schedule in every year, and 20 asset points, so that an
evaluation solves eight problems instead of thirty-six; the counterfactuals use the full size and the era
schedules.\footnote{One evaluation at this size takes 78 seconds on two threads (55 seconds for the solves) against
about 510 at the full size; five instead of seven levels would move simulated mean assets by 34\%, so seven are kept
(\code{IGE/paperA\_model/bench.py}, \code{params.SIZES}; \code{IGE/paperA\_model/README.md}).}

\paragraph{Transitory shocks in the solution.} The transitory shocks are a state on four points: within the year
the continuation value is evaluated at the current point, and in the December-to-January step it is the average over
the four points of a new draw. The work per month is therefore of the same order as with a monthly i.i.d.\ shock
integrated on nine quadrature nodes; what grows is the number of stored continuation values.

\paragraph{Storage and memory.} The solver stores only the expected continuation value $W_{t+1}(a'; \cdot)$ on the
asset grid for every state and month of the problem being solved: $211{,}680 \times 420$ values, about 0.36 GB in single
precision, plus the separated wife's problem; problems are solved and simulated one at a time, so this is the peak
(Table~\ref{tab:state}). The simulator re-solves each simulated couple's static monthly problem (the alternatives and
the saving choice) from this array, rather than storing consumption and choice-specific values for each of up to 12
alternatives. With the PI's HPC request (32 cores, \code{--threads 2}, 64 GB, 16 workers; Decision~\ref{od:hpc}) this
leaves about 4 GB per worker, within which one evaluation must fit; peak memory is measured locally before the first HPC
run.

\paragraph{Simulation.} Common random numbers: all draws (transitory, taste, displacement, job finding, human capital,
UI take-up, child, separation, couple type) are fixed across evaluations and across counterfactuals. Each simulated
couple carries its calendar month, so every moment is computed on simulated couple-months or couple-years of the same
window as its data counterpart.

\paragraph{The measurement layer.} The model produces monthly histories; each data set is reproduced by writing the
histories in that data set's shape and running the data definitions on it.
\begin{itemize}
  \item \emph{PSID}: household-years with annual hours (sum of monthly hours), annual earnings and family income (sums),
  ``displaced in year $y$'' (a displacement in any month of calendar year $y$), interview-week status (status in March),
  consumption reported at the interview (spending in the interview month, annualised), liquid wealth at the interview;
  annual waves to 1997 and biennial after; each copy observed in exactly the waves in which its PSID couple is observed.
  The Paper A scripts' definitions (\code{IGE/19}, \code{25}, \code{38}) are run on this panel.
  \item \emph{CPS}: couple-months in the four-eight-four rotation, reweighted to the CPS distribution of the wife's age by
  year; a job loss is employed in $m-1$ and not employed in $m$ after a displacement. The model has no regional or
  aggregate shocks, so its raw same-month difference is compared with the CPS estimate with state $\times$ month fixed
  effects. The added-worker effect of \citet{bacher2025joint} is built on the same panel: couples with one employed and
  one non-participating spouse, the spouse's entry into employment and into unemployment (looking) in the month the
  employed spouse loses the job, against the months in which the job continues, by the wife's age. The model's split
  between unemployment and non-participation is a measurement (Section~\ref{sec:env}), so it is compared with the data
  with and without the correction of Elsby, Hobijn and \c{S}ahin (2015) for classification error, which Bacher et al.\
  apply.
  \item \emph{SIPP}: monthly paths of wives with no job the month before the husband's involuntary separation, and of
  the husbands themselves, over 36 months; earnings in the first month of re-employment relative to the last month of
  the previous job, by spouse and length of non-employment, after an involuntary separation.
  \item \emph{DWS}: displacements by calendar year in a three-year recall window.
  \item \emph{BPS}: the biennial household-year panel with BPS's variable names and filters (Section~\ref{sec:bps}).
\end{itemize}

\paragraph{Staged checks.} Before estimation the solver passes, in order: (1) flat consumption with no risk; (2) a
precautionary buffer with risk; (3) the wife's work rises with lower assets and after the husband's loss; (4) the search
friction lengthens non-employment; (5) the wife's exits from non-employment rise when her UI entitlement lapses;
(6) with $\pi_c(\cdot) = 0$ and $m_H = 1$, the spouses' job losses are independent in the simulated CPS panel; (7) with
the transitory shocks switched off, the simulated BPS panel's transitory wage variance comes only from displacement and
human capital moves within a year. Check (3) also decides the fallback of Decision~\ref{od:search}. The coded suite
has twelve checks, 0--11 (\code{IGE/paperA\_model/tests/run\_all.py}): the endogenous grid against a brute-force
maximisation (0), checks (1)--(6) above, the joint-loss probabilities against their closed forms (7), the taste-shock
and grid profiles (8, 9), the measurement layer against the data scripts' estimators (10), and the pooled moments as
race-share averages (11); all pass at the start values. Check (7) above is run by \code{bps\_check.py}.

% #####################################################################################################
\part{Estimation and validation}

% =====================================================================================================
\section{Parameters}\label{sec:params}

All endogenous parameters are estimated jointly in one simulated-method-of-moments run; nothing is tuned by hand.
Table~\ref{tab:params} lists every parameter with its status. Calibrated parameters are fixed inputs taken from the
named source or by the stated mapping. Thirty-two parameters are estimated: the thirty of Decision~\ref{od:count},
the third common-shock probability that the workplace match adds, and the accumulation slope $\eta_{up}$
(Decision~\ref{od:eta}). At the estimation size (Section~\ref{sec:solution}) the groups $s = 1$ and $s = 2$ are merged
and share $\pi_c(1)$, so $\pi_c(2)$ does not enter and thirty-one are free (\code{params.PARAMS},
\code{estimate.py}). The code also carries the memory parameter $c_{mem}$ of the rejected alternative of
Decision~\ref{od:ridge}, held at zero. The
fixed preferences are the same as in the lifecycle model that shares this code, so that the two papers do not
contradict each other on shared objects.

\small
\begin{longtable}{L{3.2cm}L{2.2cm}L{1.6cm}L{7.1cm}}
\caption{Parameters, status and identification}\label{tab:params}\\
\toprule
Parameter & Symbol & Status & Value or source; identifying moments (estimated parameters) \\
\midrule
\endfirsthead
\toprule
Parameter & Symbol & Status & Value or source; identifying moments \\
\midrule
\endhead
\midrule \multicolumn{4}{r}{\footnotesize continued} \\
\endfoot
\bottomrule
\endlastfoot
\multicolumn{4}{l}{\textit{A. Preferences and fixed economics}} \\
Discount factor & $\beta$ & calibrated & $0.96^{1/12}$ a month (annual 0.96; Decision~\ref{od:credit}) \\
Risk aversion & $\gamma$ & calibrated & 1.5 \\
Interest rate & $r_m$, $\omega$ & calibrated & $1.03^{1/12} - 1$ on saving; borrowing wedge 8.4 points a year, the G.19 real credit-card rate 1995--2021 less 3\% (\code{borrowing\_wedge.py}; Decision~\ref{od:wedge}) \\
Equivalence scales & $s_M$, $e(k)$ & calibrated & couple 1.64 ($=1/0.61$, McClements); McClements child shares by bucket \\
Hours & $H(P)$, $H(F)$ & calibrated & 87 and 173 hours a month \\
Taste-shock scale & $\sigma_\varepsilon$ & calibrated & 0.05, not estimated; profile in Section~\ref{sec:prefs} (Decision~\ref{od:sigma}) \\
Part-time accumulation & $\kappa_{up}$ & calibrated & 0.5; profile reported \\
Consumption floor & $c_{min}$ & calibrated & AFDC/TANF maximum benefit for a family of three, by era \\
Retirement value & --- & calibrated & 240 months; Social Security 40\% of final earnings of each spouse \\
Asset division & --- & calibrated & equal at separation \\
\multicolumn{4}{l}{\textit{B. Wages and wage shocks}} \\
Wage levels & $\mu_j(r)$ & calibrated & PSID mean log hourly wage of employed wives and husbands aged 25--29, by race \\
Wage step & $\psi_j$ & calibrated & seven levels span the 1st--99th percentiles of the PSID log hourly wage residual at 40--49 \\
Persistent innovations & $p^\varepsilon_w, p^\varepsilon_h, p^c$ & calibrated & $12\psi_j^2(p^c + p^\varepsilon_j)$ equal to BPS's permanent variances (wife 0.0382, husband 0.0303), $12\psi_w\psi_h p^c$ to their covariance 0.0027 (Decision~\ref{od:wageproc}) \\
Transitory shocks & $\sigma_{\nu,w}, \sigma_{\nu,h}, \rho_\nu$ & calibrated & one draw per calendar year (Decision~\ref{od:trans}); start at BPS's transitory variances (wife 0.0125, husband 0.0275; s.d.\ 0.112, 0.166) and correlation 0.31, reduced once before estimation by the part BPS's estimator attributes to the model's other sources and the measurement-error mixture (Section~\ref{sec:labor}) \\
\multicolumn{4}{l}{\textit{C. Environment and policy inputs}} \\
Workplace match shares & $\Pr(s \mid r)$ & calibrated & copies of PSID couples, $s$ coded with the CPS definition (\code{IGE/60}); CPS shares 6.3\% ($s=2$), 11.4\% ($s=1$), 82.3\% ($s=0$) as a check (\code{IGE/37}) \\
Child process & $b_0, b_2$ & calibrated & PSID couples, monthly birth hazards by wife's age, bucket and race \\
Number of children & $n(k,\text{age},r)$ & calibrated & PSID means \\
Initial conditions & --- & data & copies of PSID couples (Section~\ref{sec:init}) \\
UI & $\rho_{UI}$, $\bar b_e$, expiry & calibrated & replacement 0.5; mean state maximum weekly benefit by era (\code{IGE/12b.UIGenerosity.py}); expiry probability $1/6$ a month (26 weeks expected), as \citet{bacher2025joint} \\
UI take-up & $q_{UI}$ & calibrated & share of displaced workers receiving UI, Displaced Worker Supplement (to compute in \code{IGE/39}) \\
EITC & $\epsilon^{in}, \bar E,$ $\epsilon^{out}, Y^{out}$ & calibrated & IRS / Tax Policy Center parameters by year and number of children \\
SNAP & $A, D$, limits & calibrated & USDA FNS allotments and deductions by fiscal year; Census poverty lines \\
Income and payroll tax & $\chi_e, \mu_e$ & calibrated & log-linear fit of federal and payroll taxes of joint filers by era (TAXSIM), BPS's form \\
Childcare price & $p_{cc}(k)$ & calibrated & PSID childcare spending per hour of the mother's work by child age \\
\multicolumn{4}{l}{\textit{D. Displacement risk (estimated)}} \\
Wife's own hazard & $\delta_w(B), \delta_w(W)$ & estimated & her monthly non-temporary job-loss rate when he keeps his job (CPS), by race; her annual displacement rate (PSID) \\
Husband's own hazard & $\delta_h(B), \delta_h(W)$ & estimated & his annual displacement rate, 8.0\% (Black) and 4.8\% (White) (PSID) \\
Common shock & $\pi_c(0), \pi_c(1), \pi_c(2)$ & estimated & same-month excess by workplace match (CPS); same-year joint displacement by sector and three-digit industry (PSID); at the estimation size $s = 1, 2$ share $\pi_c(1)$ and $\pi_c(2)$ does not enter \\
Couple type & $p_H, m_H$ & estimated & the wife's excess in the three months before and after his loss (CPS); permanent component and placebo window (PSID); start on the CPS ridge at $p_H = 0.02$, $m_H = 15.95$ (Decision~\ref{od:ridge}) \\
\multicolumn{4}{l}{\textit{E. Search friction (estimated)}} \\
Wife's job finding & $\lambda_w(B), \lambda_w(W)$ & estimated & persistent monthly employment entry of non-employed wives (CPS), by race (Decision~\ref{od:persist}); her employment path after his separation, months 0--6 (SIPP) \\
Husband's job finding & $\lambda_h(B), \lambda_h(W)$ & estimated jointly (Decision~\ref{od:husband}) & his employment 3, 6 and 12 months after an involuntary separation (SIPP); his hours in years 0--1 after displacement (PSID), by race; job finding by age validated \\
\multicolumn{4}{l}{\textit{F. Human capital (estimated)}} \\
Wife's accumulation & $\pi^{up}_w$ & estimated & two-year log wage growth of continuously employed wives, by age group (PSID) \\
Wife's depreciation & $\pi^{dn}_w$ & estimated & her earnings on re-employment by length of non-employment after an involuntary separation (SIPP); log wage at re-entry after 12+ months out relative to continuously employed wives (PSID) \\
Wife's displacement loss & $d_w$ & estimated & her labor income loss after her own displacement, about \$4{,}600 a year (PSID); her loss on re-employment after 1--3 months (SIPP); cut if that loss is close to zero \\
Husband's accumulation & $\pi^{up}_h$ & estimated & two-year log wage growth of continuously employed husbands, by age group (PSID) \\
Husband's depreciation & $\pi^{dn}_h$ & estimated & his earnings on re-employment by length of non-employment (SIPP); his labor income in year 6 relative to years 1--3 after displacement (PSID) \\
Husband's displacement loss & $d_h$ & estimated & his labor income loss in years 1--3, $-17\%$ / $-24\%$ (PSID); his loss on re-employment after 1--3 months (SIPP); cut if that loss is close to zero \\
Accumulation by level & $\eta_{up}$ & estimated (Decision~\ref{od:eta}) & the fall of the two-year wage growth of the continuously employed at 35--44 and 45--54 relative to 25--34, wives and husbands (PSID); start 0.5, bounds $[0, 3]$ \\
\multicolumn{4}{l}{\textit{G. Work, safety net, marriage, credit (estimated)}} \\
Wife's full-time cost & $\phi_F$ & estimated & employment of wives with no child under 18 (PSID); wives' persistent monthly exit from employment outside job loss (CPS) \\
Wife's part-time cost & $\phi_P$ & estimated & part-time share of employed wives (PSID) \\
Child cost, 0--5 & $\theta_{05}$ & estimated & employment of wives with youngest child 0--5; her added hours and employment after his displacement, youngest 0--5 (PSID) \\
Child cost, 6--17 & $\theta_{617}$ & estimated & the same with youngest child 6--17 (PSID) \\
Husband's full-time cost & $\phi_h$ & estimated & part-time share of employed husbands (PSID) \\
Fixed cost of work & $F$ & estimated & consumption change after her own displacement (PSID 1999--2021); joint-minus-single consumption change (PSID 1999--2021); employment-consumption complementarity (validation) \\
SNAP stigma & $\kappa_S$ & estimated & SNAP receipt of couples; its rise after his displacement, 4.2 / 3.3 points (PSID) \\
Baseline separation & $\sigma_0(B), \sigma_0(W)$ & estimated & annual separation rate of couples, by race (PSID) \\
Separation response & $\sigma_J(B), \sigma_J(W)$ & estimated & still with wife, $-4.5$ / $-1.9$ points in years 1--3 (PSID) \\
Borrowing limit & $\underline a$ & estimated (Decision~\ref{od:credit}) & share of couples with liquid wealth $\le 0$ by wife's age group; consumption change after his displacement (PSID 1999--2021) \\
\end{longtable}
\normalsize

% =====================================================================================================
\section{Targets}\label{sec:targets}

Table~\ref{tab:targets} lists the moments. Each moment is computed on simulated households of the same window as its
data counterpart: the displacement and hours moments on couple-years of 1984--2021, the consumption moments on
couple-years of 1999--2021 only, the CPS moments on couple-months of 1994--2023 and the SIPP moments on 2003--2016. The
consumption window is stated separately because PSID expenditure covers only 1999--2021 (Decision~\ref{od:cons}).
Standard errors are bootstrap standard errors (200 draws of 1968 families in the PSID, of households in the CPS and
SIPP) produced by one targets script (\code{IGE/60.PaperATargets.py}), which also writes
\code{EstimationMoments/couples\_*.csv}; the consumption moments come from \code{IGE/47.PSIDConsumptionEvent.py}
(done), with standard errors clustered by 1968 family.
Where Paper A reports the moment on a different window or definition, the table says so and the targets script
recomputes it on the stated window.

\paragraph{The CPS flows (Decision~\ref{od:persist}).} The wives' entry and exit rates are persistent flows, computed
with one filter in the data (\code{IGE/60}) and in the simulated CPS panel (\code{panels.cps\_samples}): a transition
counts only if the old status held in months $t-1$ and $t$ and the new one holds in $t+1$ and $t+2$, all four months in
one rotation spell, with the husband employed in $t$ and $t+1$. The filter removes one-month reversals, which the model's
monthly statuses do not produce; it cuts the data's exit rate from 0.0261 to 0.0101 and the entry rates from 0.0751 to
0.0320 (White) and from 0.0942 to 0.0404 (Black), while the model's flows are almost all persistent
already.\footnote{At the extended start the model's one-month and persistent rates are 0.0023 and 0.0014 (exit), 0.047
and 0.036 (White entry), 0.099 and 0.089 (Black entry); with the persistent flows the objective at the start falls from
90{,}296 to 12{,}761. The response-error correction of \citet{poterba1986reporting} was coded first and dropped
(\code{estimate.py --cps-persist off --cps-classif on}; error rates for women 25--59 in
\code{moments.CPS\_CLASSIF\_PS86\_W2559}): it brings the simulated exit rate to 0.0240 against 0.0261, but the entry
rates to 0.208 and 0.364 against 0.075 and 0.094; at the data's employment rate of 0.81 the spurious entries of
misrecorded employed wives alone give a measured entry rate of at least 0.083, above the White target, so the true entry
rate their errors imply is negative. Their errors are for 1981 and serially uncorrelated, which they note overstates the
adjustment, and the CPS targets start in 1994, after the redesign with dependent interviewing. Code:
\code{IGE/60.PaperATargets.py}, \code{IGE/paperA\_model/panels.py}; \code{results/diagnostics\_20261018.md}, step 1 and
follow-up 2.} The persistent exit rate is still the largest misfit at the start: 0.0014 against 0.0101, 4{,}623 of an
objective of 11{,}734 (11{,}733.7 from a clean compilation on the PC, 11{,}733.6 on the cluster; compiled binaries
differ in the fourth significant digit, so the reproducibility check compares the objective at a relative $10^{-3}$).

\paragraph{The objective.} Of the 97 targets, 85 enter the objective (\code{moments.OBJ\_EXCLUDE}). The other twelve stay
in the fit table, one reason each: the four pooled CPS rates (in the model each is exactly the race-share average of its
Black and White rates; the data's also contain other races), the three one-month CPS flows (replaced by the persistent
ones), the three SIPP employment rates of husbands without a separation (in the model husbands never quit), and the
husband's employment rate by race (Decision~\ref{od:hemp}). The weights are $1/(\mathrm{se}_m^2 +
\mathrm{var}^{sim}_m)$: the bootstrap variance of each data moment plus the simulation variance of the model moment,
from five sets of common random numbers (\code{inputs/moment\_noise.csv}, scaled to the run's number of copies), so that
a moment the simulation reproduces noisily does not dominate; block weighting is a robustness run.\footnote{Under
inverse bootstrap variances the CPS hazards held 95.6\% of the diagonal weight at the start; with the trim and the
simulation variance 90.6\%, the wife's exits 8.1\%. The simulation variances are those of the start before the
extensions: at the extended start the one-month entry rates corrected for response error had a simulation standard
deviation seven times larger (0.0004 to 0.0030), which would have removed the search block from the objective. They
are recomputed at the Levenberg--Marquardt candidate (job \code{Joint-noise}) before the global search
(\code{results/diagnostics\_20261018.md}).}

\small
\begin{longtable}{L{1.9cm}L{4.6cm}L{2.8cm}L{1.6cm}L{2.9cm}}
\caption{Targets, windows and sources}\label{tab:targets}\\
\toprule
Block & Moment & Value (s.e.) & Window & Source; status \\
\midrule
\endfirsthead
\toprule
Block & Moment & Value (s.e.) & Window & Source; status \\
\midrule
\endhead
\midrule \multicolumn{5}{r}{\footnotesize continued} \\
\endfoot
\bottomrule
\endlastfoot
Joint loss, monthly & Wife's same-month job-loss excess when he loses his, non-temporary losses (state $\times$ month FE) & 0.026$^{***}$ on a base of 0.0023 (headline, all losses: 0.040, s.e.\ 0.002) & CPS 1994--2023 & \code{IGE/37}; s.e.\ to report \\
 & Same, by workplace match: shared-employer cell / same detailed industry / outside the shared-employer cell & 0.337$^{***}$ (0.015) / 0.211$^{***}$ (0.009) / 0.026$^{***}$ (0.001) (all losses) & CPS 1994--2023 & \code{IGE/37}; recompute non-temporary, three exclusive groups, in \code{IGE/60} \\
 & Wife's excess in months $-3, -2, -1$ relative to his loss & 0.006, 0.006, 0.007 (all $p<0.001$) & CPS 1994--2023 & \code{IGE/30}; recompute non-temporary \\
 & Wife's excess in months $+1, +2, +3$ & 0.008, 0.006, 0.008 (all $p<0.001$) & CPS 1994--2023 & \code{IGE/30}; recompute non-temporary \\
Joint loss, annual & Permanent couple component & 0.020$^{**}$ (0.008) & PSID 1987--2021 & \code{IGE/38}; s.e.\ implied by $p=0.012$ \\
 & Contemporaneous component & 0.033$^{**}$ (0.016) & PSID 1987--2021 & \code{IGE/38}; s.e.\ implied by $p=0.036$ \\
 & P(she displaced $\mid$ he is) $-$ P($\mid$ he is not), same interval & 0.102 $-$ 0.039 $=$ 0.063$^{***}$ (0.005) & PSID 1984--2021 & \code{IGE/25} \\
 & Same, same sector & 0.211 $-$ 0.028 $=$ 0.183$^{***}$ (0.022) & PSID 1984--2021 & \code{IGE/25} \\
 & Same, different sector & 0.076 $-$ 0.035 $=$ 0.041$^{***}$ (0.006) & PSID 1984--2021 & \code{IGE/25} \\
Base hazards & Wife's monthly non-temporary job-loss rate, he employed, by race & B 0.0034 (0.0001), W 0.0022 (0.00004); pooled 0.0023 in the fit table only & CPS 1994--2023 & \code{IGE/37} (\code{IGE/60}) \\
 & Husband's annual displacement rate, working male heads & B 0.080, W 0.048 & PSID 1984--2021 & \code{IGE/19}; recompute on window \\
Her search & Persistent monthly employment entry of non-employed wives, husband employed (Decision~\ref{od:persist}) & B 0.0404 (0.0019), W 0.0320 (0.0004); one-month B 0.094, W 0.075 and the pooled rates in the fit table only & CPS 1994--2023 & \code{IGE/60} \\
 & Employment of wives with no job the month before, relative to wives whose husband stays employed, months 0, 1, 6 after his involuntary separation & 0.019$^{***}$, 0.037$^{***}$, 0.054$^{***}$ & SIPP 2003--2016 & \code{IGE/40} \\
Her exits & Wives' persistent monthly exit from employment (to unemployment or non-participation) without a job-loss reason, husband employed & 0.0101 (0.0001); one-month 0.0261 in the fit table only & CPS 1994--2023 & \code{IGE/60} \\
His search & His employment 3, 6, 12 months after an involuntary separation (vs.\ no separation) & 0.586, 0.708, 0.791 (0.983, 0.972, 0.960; the no-separation rows in the fit table only) & SIPP 2003--2016 & \code{IGE/40} (\code{IGE/60}) \\
His employment & Husband employed (positive hours) in the couple-year, by race & B 0.882 (0.005), W 0.943 (0.002); fit table only (Decision~\ref{od:hemp}) & PSID 1984--2021 & \code{IGE/60} \\
 & His annual hours, years 1--3 after displacement & B $-146^{***}$ (35), W $-130^{***}$ (28) & PSID 1984--2021 & \code{IGE/19}; recompute on window \\
His earnings & Labor income, years 1--3 & B $-4{,}160^{***}$ (757), W $-9{,}622^{***}$ (1{,}071) & PSID 1984--2021 & \code{IGE/19}; recompute on window \\
 & Labor income, year 6 & about $-6{,}000$ to $-7{,}000$ & PSID 1984--2021 & \code{IGE/19}; s.e.\ to report \\
Her earnings & Labor income after her own displacement, years 1--3 & about $-4{,}600$ & PSID 1984--2021 & \code{IGE/25}; s.e.\ to report \\
Wage growth & Two-year log wage growth, continuously employed wives and husbands, ages 25--34, 35--44, 45--54 & to compute & PSID 1984--2021 & \code{IGE/46} \\
 & Her log wage at re-entry after 12+ months out, relative to continuously employed & to compute & PSID 1984--2021 & \code{IGE/46} \\
 & Earnings in the first month of re-employment relative to the last month of the previous job, after 1--3, 4--12, 13--24 months of non-employment, involuntary separations, wives and husbands & to compute & SIPP 2003--2016 & \code{IGE/40}; in \code{IGE/60} \\
Added work & Wife's hours, years 1--3, youngest 0--5 / 6--17 & $+112^{***}$ (31) / $-26$ (36) & PSID 1984--2021 & \code{IGE/19}; published on all events, recompute from 1984 \\
 & Wife employed, years 1--3, youngest 0--5 / 6--17 & $+0.041^{**}$ (0.018) / $-0.038^{**}$ (0.018) & PSID 1984--2021 & \code{IGE/19}; recompute from 1984 \\
 & Gross response (wife kept her job), all / youngest 0--5 & $+47$ ($p=0.165$) / $+97^{**}$ ($p=0.032$) & PSID 1984--2021 & \code{IGE/38} \\
 & Offset from joint losses & $-35^{**}$ (14) & PSID 1984--2021 & \code{IGE/38} \\
Employment levels & Wives' employment by youngest child (none, 0--5, 6--17); part-time share of employed wives; part-time share of employed husbands & to compute & PSID 1984--2021 & \code{IGE/46} \\
Safety net & SNAP receipt after his displacement, years 1--3 & B $+0.042^{**}$ (0.019), W $+0.033^{***}$ (0.010) & PSID 1984--2021 & \code{IGE/19}; recompute on window \\
 & SNAP receipt of couples, by race & to compute & PSID 1984--2021 & \code{IGE/46} \\
Separation & Still with wife, years 1--3 & B $-0.045^{***}$ (0.015), W $-0.019^{**}$ (0.009) & PSID 1984--2021 & \code{IGE/19}; recompute on window \\
 & Annual separation rate of couples, by race & to compute & PSID 1984--2021 & \code{IGE/46} \\
Wealth & Share of couples with liquid wealth $\le 0$ and above one year of earnings, by wife's age group (25--34, 35--44, 45--59) & to compute & PSID waves 1984, 1989, 1994, 1999--2021 & \code{ReducedForm/17} for couples \\
Consumption & Change in log total consumption, first three interviews after his single displacement, married couples: pooled; Black / White & $-0.013$ (0.019; $p=0.492$); $-0.005$ (0.031; $p=0.863$) / $-0.005$ (0.023; $p=0.820$) & PSID 1999--2021 (events 2000--2020) & \code{IGE/47.PSIDConsumptionEvent.py} \\
 & Same after her own displacement (he not displaced) & $-0.031$ (0.024; $p=0.200$) & PSID 1999--2021 & \code{IGE/47.PSIDConsumptionEvent.py} \\
 & Joint minus single displacement (63 joint events; target because s.e.\ $< 0.08$) & $-0.059$ (0.073; $p=0.422$) & PSID 1999--2021 & \code{IGE/47.PSIDConsumptionEvent.py} \\
\end{longtable}
\normalsize

\noindent The consumption measure is BPS's (nondurables and services with a 6\% rent equivalent), ported from their
\code{psid\_consumption.do} to Python; event time is in waves, dated as in the PSID (a displacement in the calendar year
the job ended; consumption reported at the next interview), and the same dating is applied to the simulated panel.
Measurement error in consumption with BPS's variance ($E(\Delta c_t\Delta c_{t-2}) = -0.0332$) is added to simulated
consumption before the moment is computed.

% =====================================================================================================
\section{Identification}\label{sec:ident}

\paragraph{The three parts of displacement risk.} The common shock $\pi_c(s)$ produces excess joint loss only in the
month it hits: the CPS spike in his loss month in each of the three workplace groups, with the PSID same-year moments by
sector over-identifying it. The couple type $(p_H, m_H)$ produces a flat elevation of the wife's hazard in every month around
the husband's loss (the CPS leads and lags) and the PSID placebo-window excess (the permanent component); because it
raises both own hazards for years, it is also what makes the risk matter for saving. The share and the multiplier are
separated by the level of the leads and lags relative to the base rates (a few couples with much higher risk against
many with slightly higher risk change the base rates differently from the excess). The own hazards $\delta_j(r)$ are
pinned by the base rates almost one for one. In an annual model the common and permanent parts collapse into one
same-year correlation; the monthly period is what separates them.

\paragraph{Search friction and work costs.} The wife's job finding $\lambda_w(r)$ is identified by the monthly entry
rate of non-employed wives and by the speed with which employment builds up after the husband's loss, while the work
costs are identified by the levels: employment by child age ($\phi_F$, $\theta_{05}$, $\theta_{617}$), the part-time
share ($\phi_P$). The response to his loss separates the two: a high work cost lowers the eventual response, a high search
friction delays it. The wives' persistent monthly exit from employment outside job loss was to pin, through the same
work costs, the margin on which a wife leaves a job; at the start values the work costs alone miss it sevenfold, and the
evidence on why wives leave employment decides what the model needs for it. The husband's $\lambda_h(r)$ is identified by his
re-employment path.

\paragraph{Human capital.} Accumulation is identified by wage growth of the continuously employed: its level at 25--34
by $\pi^{up}_j$, its fall with age by $\eta_{up}$. The
earnings on re-employment by the length of non-employment separate depreciation from the displacement loss:
depreciation sets the slope of the loss in the length of the spell, $d_j$ its level after a short spell. The earnings
losses after displacement (with hours controlled by the hours path), the wage at re-entry (wife) and the persistence of
his loss from years 1--3 to year 6 (husband) over-identify both.

\paragraph{Safety net, marriage and credit.} The SNAP stigma $\kappa_S$ is identified by receipt among couples and its
rise after displacement, the separation parameters by the baseline rates and the response, the borrowing limit by the
share of couples with no liquid wealth and by the consumption fall after displacement, and the fixed cost $F$ by the
consumption change at the wife's own displacement (spending falls less than income because the cost of work ends).

\paragraph{What the added-worker effect by age tests.} No parameter depends on age: the work costs depend on the
youngest child's age, the search friction on race and spouse, the hazards on race, workplace match and couple type.
An age gradient in the simulated added-worker effect of one-earner couples can therefore come only from the states that
change with age---both spouses' human capital, assets, the children---and from the approach of the terminal value at 60.
\citet{bacher2025joint} attribute the decline of the effect from 7.1 points at 26--35 to 1.4 points at 56--65 to the
first two and to the horizon. Because the effect is not targeted (Decision~\ref{od:awe}), matching or missing its
gradient is a test of those channels; if the simulated effect at 46--55 is far from the data, the terminal value (240
months, Social Security at 40\% of final earnings) is the first suspect.

\paragraph{What the correlation buys.} A restricted model with independent shocks ($\pi_c(\cdot) = 0$, $m_H = 1$) is
estimated on the same targets with the same draws and reported next to the main model, in the fit table and in the
counterfactual table. The difference in fit and in counterfactuals is what the correlation buys.

\paragraph{Checks and procedure.} A Jacobian of the moments at the start values checks that every parameter moves a
named block. At the extended start (31 free parameters) the range-scaled condition number is 529 and the weakest
directions are $m_H$ with $\eta_{up}$, $\pi^{dn}_h$ with $d_h$, $\lambda_h(B)$ against $\lambda_h(W)$, and $F$; the
Jacobian is read again at the candidate, and ties across races remain an option of the code (\code{--tie}), not
applied. The search runs in stages on the same common random numbers ($K = 3$ copies of each PSID couple):
closed-form start values (\code{analytic\_start.py}); a Levenberg--Marquardt search in logistic units with a
central-difference Jacobian at each iteration (\code{estimate.py polish}); the simulation variances recomputed at its
candidate (\code{estimate.py noise}); CMA-ES around the candidate, with initial step sizes from the Jacobian's column
norms and three seeds as the check against local minima (\code{estimate.py run --cma-stds-from}); sandwich standard
errors with the bootstrap covariance of the targets and the simulation covariance of the moments
(\code{estimate.py se}). The same route is first run on pseudo-data generated at known parameters (\code{polish
--pseudo}), and the independent-shock model is estimated by it.\footnote{CMA-ES alone was tried first: a pseudo-data
recovery of 19 parameters from 496 evaluations lowered the objective from 2{,}907 to 482 (zero at the truth) and was
still falling, with 3 of 19 parameters within two simulation standard errors of the truth and 8 within two data
standard errors (\code{results/runs/recovery\_closed\_20261018\_151959/recovery\_report.md}).}

% =====================================================================================================
\section{Untargeted validation}\label{sec:valid}

The following moments are computed on the simulated households through the same definitions and windows and are not
used in estimation.
\begin{enumerate}
  \item \textbf{CPS search jump}: non-employed wives unemployed and looking $+8.8$ points (0.088$^{***}$, s.e.\ 0.004) the
  month after his loss, on a base of 6.0\%; Black 0.103$^{***}$, White 0.083$^{***}$; and the CPS employment path of
  these wives, $+0.022^{***}$, $+0.033^{***}$, $+0.042^{***}$ and $+0.040^{***}$ after 1, 2, 3 and 12 months
  (\code{IGE/30}).
  \item \textbf{SIPP 36-month path}: among wives with no job the month before, looking $+0.053^{***}$ at once, fading to
  0.013 ($p=0.246$) at 36 months; employment 0.044$^{***}$ after 24 months; net change in all wives' employment about zero
  for a year ($-0.007^{**}$ after one month, 0.006, $p=0.297$, after twelve) because working wives lose their own jobs
  (\code{IGE/40}). Months 0--6 are targeted; months 12--36 are validation.
  \item \textbf{DWS same-year ratio}: a wife displaced in the same calendar year as her husband in 4.7\% of couples in which he
  was displaced against 1.1\% (ratio 4.3; excess 0.033$^{***}$, s.e.\ 0.001); tetrachoric correlation 0.26 (\code{IGE/39}).
  \item \textbf{BPS on simulated data}: Section~\ref{sec:bps}.
  \item \textbf{Food and consumption}: food spending after single and joint displacement in dual-earner couples
  ($-\$419$, $p=0.092$, and $-\$699$, $p=0.091$; difference $-0.047$ in logs, $p=0.639$; \code{IGE/38}); in BPS's panel,
  consumption after a joint relative to a single displacement, about zero in the same two-year interval and $-0.07$ to
  $-0.14$ in the next ($-0.071$, $p=0.19$, BPS sample; $-0.138^{**}$, $p=0.027$, with the low-income sample), checked as a
  sign and a band on the simulated BPS panel with BPS's residualisation (1999--2009; \code{IGE/bps\_reconciliation/}).
  \item \textbf{Family income} after his displacement, years 1--3: B $-4{,}832^{***}$ (1{,}082), W $-11{,}179^{***}$
  (1{,}353) (\code{IGE/19}).
  \item \textbf{Wives already unemployed} find work less often after his loss ($-0.020^{*}$, $-0.050^{**}$ after one and
  three months; \code{IGE/30}); the model, in which the common shock does not lower her job finding, is expected to miss
  this, and the miss is reported.
  \item \textbf{The excess by child age and by race}: equal in the data (CPS; PSID same-year 0.065 youngest 0--5, 0.050
  youngest 6--17, 0.072 no child).
  \item \textbf{Assets relative to human wealth}: $\pi(\text{age}) = a/(a + \text{human wealth})$ against BPS's Figure 3
  (0.03 at 30, 0.15 on average, 0.36 after 55).
  \item \textbf{Husbands' earnings by age}: the level of the PSID married men's profile
  (\code{ModelInputs/}\allowbreak\code{male\_labinc\_age\_by\_}\allowbreak\code{race\_edu\_smoothed.csv}).
  \item \textbf{Added-worker effect by age} (Decision~\ref{od:awe}): the definition of \citet{bacher2025joint} on Paper
  A's CPS panel---couples with one employed and one non-participating spouse, the spouse's entry in the month the
  employed spouse loses the job---overall, by wife's age 26--35, 36--45, 46--55 and 56--60 (Decision~\ref{od:age}), and
  split into entry into employment and into unemployment. Their estimates, 5.55 points overall and 7.12, 6.83, 4.57 and
  1.40 points from the youngest to the oldest group (56--65), are the reference (\code{IGE/50.CPSAddedWorkerByAge.py}, in
  progress).
  \item \textbf{The joint-loss excess by age}: the wife's same-month job-loss excess (state $\times$ month fixed effects,
  all and non-temporary losses) by the same age groups (\code{IGE/50}). The model's own hazards and common shock do not
  depend on age, so an age pattern in the data would show what the model leaves out.
  \item \textbf{Joint labor-market states by age}: the shares of couples with both employed, one employed and one out of
  the labor force, one employed and one unemployed, and both unemployed, by wife's age (the data counterpart of
  \citealp{bacher2025joint}, their Figure 5; \code{IGE/50}). It checks that the children state does the work their
  age-dependent disutility does; the both-unemployed share is the stock counterpart of the joint-loss flow.
  \item \textbf{Transitions}: the monthly employment--unemployment--non-participation transition matrix of wives and of
  husbands (CPS, \code{IGE/30}); unemployment and non-participation are a measurement in the model.
  \item \textbf{Job finding by age}: the simulated monthly job finding of non-employed wives and husbands by age group,
  against the CPS and SIPP rates by age and against the age profile in \citet{bacher2025joint} (their Figure 10;
  Decision~\ref{od:husband}).
\end{enumerate}

% =====================================================================================================
\section{BPS on simulated data}\label{sec:bps}

The test asks whether one model, estimated on the joint-loss moments, gives BPS's estimates when their estimator runs on
a BPS-style sample of its own simulated households, and the joint-loss facts on the full sample. It is run once, at the
estimated parameters, and is not part of estimation.

\subsection{Protocol}
\begin{enumerate}
  \item \textbf{Simulate} about 50{,}000 couples (copies of PSID SRC couples, the non-SEO sample BPS use), monthly, with
  calendar years labelled 1998--2008 so that BPS's hard-coded tax schedules and year dummies apply.
  \item \textbf{Aggregate} to a PSID-shaped household-year panel with BPS's variable names: \code{ly}, \code{wly}
  (annual earnings), \code{hours}, \code{hourw} (annual hours), \code{totcons} (annual consumption), \code{tot\_assets3}
  (start-of-year assets), \code{empst}, \code{wempst} (status in a spring month: working, temporarily laid off,
  looking), \code{fsize}, \code{kids}, \code{educ}, \code{weduc}, \code{race}, \code{state}, \code{marit}, \code{age},
  \code{agew}, \code{fchg} (2 or more in a separation year), \code{year} in BPS's coding, \code{seo} $=0$,
  \code{mortgage1}, \code{mortgage2}, \code{tyoth}, \code{outkid}, \code{smsa} (\code{sim\_to\_bps.py}).
  \item \textbf{Add measurement error} as a two-component mixture on earnings, hours and consumption: a common small error
  with Bound et al.'s shares of the log variance (earnings 4\%, hours 23\%, wages 13\%; \citealp{bound1994evidence}; BPS's
  \code{T4C1.do}) and a rare gross error affecting a few percent of observations, large enough that BPS's trims select
  mostly gross errors; plus consumption error with variance matching $E(\Delta c_t \Delta c_{t-2}) = -0.0332$. The mixture
  is calibrated to two diagnostics of the PSID: the composition of the trimmed rows (of 113 rows removed by the
  half-minimum-wage rule, 7 are husband-displacement, 5 wife-displacement and 0 joint-displacement rows; of 41 rows removed
  by the jump trim, 0, 1 and 0) and the ratio of the husband's transitory wage variance without and with the jump trim
  (0.056 / 0.028).\footnote{Without the mixture, the 0.25-percentile tail of $\Delta\log x_t \times \Delta\log x_{t-2}$ in
  simulated data would consist of genuine collapse-and-recovery years, the displacements, and BPS's trim would remove
  exactly the events the test is about; in the PSID the trims remove measurement-error outliers, not displacement years.
  Code: \code{IGE/bps\_reconciliation/} (Table S2.2).}
  \item \textbf{Run BPS's code unchanged} on odd waves only (biennial): \code{psid\_sample.do} (from line 80),
  \code{prepare\_bootstrap.do}, \code{residual\_measures.do}, and \code{bootstrap.do} with \code{bn = 0} (which calls
  \code{tax\_rate\_estimation\_snap.do}, \code{pi\_s\_estimation.do}, \code{variance\_estimation.do},
  \code{frisch\_estimation.do}), the insurance decomposition and the Table 7 regressions, through
  \code{IGE/bps\_reconciliation/code/bps\_on\_sim.do}, which is first tested on BPS's own panel (it must reproduce the
  replication exactly).
  \item \textbf{Relax each restriction} on the simulated panel as on the real one (the year-level instead of person-level
  no-work drop; no status controls; trims off; low-income couples in; all), and run the Step 2 displacement event study.
  \item \textbf{On the full monthly sample}, compute Paper A's moments through the measurement layer of
  Section~\ref{sec:solution}.
\end{enumerate}

\subsection{Pass criteria}

Table~\ref{tab:bpspass} states the criteria in advance. Not expected to match: spousal leisure complementarity (not
modelled) and exact Frisch magnitudes (discrete hours give the signs and the decomposition, not the elasticities).

\small
\begin{longtable}{L{6.4cm}L{8.4cm}}
\caption{Pass criteria for BPS on simulated data}\label{tab:bpspass}\\
\toprule
BPS estimate or PSID derivative & The simulated BPS-style panel must give \\
\midrule
\endfirsthead
\toprule
BPS estimate or PSID derivative & The simulated BPS-style panel must give \\
\midrule
\endhead
\bottomrule
\endlastfoot
\multicolumn{2}{l}{\textit{A. Levels (BPS baseline, replicated)}} \\
Wage variances and covariances (Table 3): permanent 0.0303 / 0.0382, cov.\ 0.0027; transitory 0.0275 / 0.0125, cov.\ 0.0058$^{**}$ & within BPS's standard errors \\
Insurance split 62.3 / 19.4 / 18.3\% (family / government / savings) & each share within about 10 points \\
Outside insurance $\beta = -0.138$ (s.e.\ 0.270) & within its standard error \\
Extensive margin $\theta_1 = -0.168^{*}$ (s.e.\ 0.087) & within its standard error \\
Own-wage and consumption-wage elasticities (Table 4--5) & same signs \\
\multicolumn{2}{l}{\textit{B. Derivatives with respect to BPS's restrictions (PSID, Steps 1--2)}} \\
No-work rule at the year level: family share 0.62 $\to$ 0.56, government 0.19 $\to$ 0.22, savings 0.18 $\to$ 0.22 & same direction, roughly the same size (the husband's long spells entering the moments) \\
Status controls removed: accounting unchanged; wife's earnings residual at her displacement doubles ($-0.101$ to $-0.202^{***}$) & the same \\
Low-income couples in: $E(\pi)$ 0.15 $\to$ 0.14, family share up (0.68) & same direction when the model's low-asset couples are added \\
Trims off: transitory variance doubles & reproduced through the measurement-error mixture \\
Attrition: of 57 joint displacements, 28 reach the sample and 17 carry a wage moment; the no-work rule drops about 10\% of male-head observations & similar shares through the same filters (the 10\% requires his spells to span a calendar year: a check on $\lambda_h$) \\
P(wife displaced in the interval $\mid$ husband displaced) 6.8\% vs 2.9\% (BPS panel); 10.2\% vs 3.9\% per couple-year (Paper A) & both, from the same shock process \\
Joint displacement years raise $E(\Delta u_{y_1}\Delta u_{y_2})$ by $+0.19^{**}$ and the wage covariance by nothing & the same: the joint loss moves hours and earnings; the hourly wage moves only through the displacement loss \\
$\theta_1$ is entries ($-0.159^{**}$), not exits ($-0.009$) & the same split \\
\multicolumn{2}{l}{\textit{C. The central reconciliation number}} \\
Wife's Marshallian hours response to the husband's permanent wage, $-0.75^{***}$ on BPS's sample & about $-0.75$ on the simulated BPS sample, and $+50$ to $+100$ hours after a displacement on the full sample \\
\end{longtable}
\normalsize

\noindent The BPS test runs on the pooled non-SEO-type sample, as BPS do; the White-only simulated subsample is compared
with the reconciliation's White column (0.52 / 0.21 / 0.27). If the model fails on a named moment, the failure is reported
as informative (for example, a family share that the model cannot produce without BPS's trims).

% =====================================================================================================
\section{Counterfactuals}\label{sec:cf}

Each experiment re-solves the model with the same draws, is reported for the main and the independent-shock models,
and is scored on: consumption in the month of the husband's loss and over the following 12 months; the share of
couple-months with no earner and the share with no earner and no credit; a consumption-equivalent welfare change at
entry; the children's adult outcomes under the high and low slopes of Section~\ref{sec:children}; each by workplace
match and race. The welfare change covers the married household and, after a separation, the wife's household; the
husband is not followed after a separation (Decision~\ref{od:welfare}), so the welfare numbers omit his part of the
couple's loss from a separation, a limitation stated with every welfare number.
\begin{enumerate}
  \item \textbf{Common shock off}: $\pi_c(\cdot) = 0$. Measures how much consumption smoothing and how much of the
  children's outcomes the correlation costs, by workplace match.
  \item \textbf{UI for both spouses and longer in recessions}: (a) at a joint loss both spouses are entitled ($q_{UI} = 1$
  for both); (b) duration from 26 to 52 weeks (expiry probability from $1/6$ to $1/12$ a month), evaluated at the
  estimated $\pi_c$ and with $\pi_c$ doubled (the recession excess, 6.7 against 3.6 points), as a comparative static
  without an aggregate state. The added-worker effect of one-earner couples is reported under (b) next to
  \citet{bacher2025joint}, their Table 12: extending expected benefit duration from six to twelve months lowers their
  model's effect from 5.16 to 3.68 points (young couples 9.35 to 6.93, old 1.22 to 0.81). The experiment here differs in
  that both spouses can be entitled in the same month after a common shock, and in the version with $\pi_c$ doubled.
  \item \textbf{Earnings subsidy against an unconditional transfer of equal cost}: an EITC-type expansion against a
  monthly transfer to couples with children whose simulated cost is equal; scored on consumption in joint-loss months, on
  months with no earner and no credit, and on children. The earnings subsidy pays only when someone works, so it insures
  least in the state in which both spouses have lost their jobs. Before this counterfactual is written, Haan and Prowse
  (2024, \emph{American Economic Journal: Macroeconomics}), who compare social assistance with unemployment insurance in a
  life-cycle model of family labor supply, are to be read and the design compared with theirs.
  \item \textbf{Childcare subsidy} when the youngest child is under six ($p_{cc}(1)$ halved and set to zero): her
  employment, her added work after his loss, children through income.
  \item \textbf{Black couples with White displacement hazards} ($\delta_w$, $\delta_h$ and the couple-type
  distribution set to the White values), \textbf{then with the White separation response} ($\sigma_J$), one at a time.
\end{enumerate}

% =====================================================================================================
\section{Code and compute}\label{sec:build}

\paragraph{The code.} The model is coded in one self-contained folder, \code{IGE/paperA\_model/}, which reads the
targets script and the BPS panel writer of the reconciliation without changing them: \code{params.py} (the parameters
with start values and bounds, fixed inputs, era schedules, solver sizes), \code{analytic\_start.py} (closed-form start
values and the wage-process mapping), \code{model.py} and \code{solve.py} (one problem $(r, s, \tau)$ in one era,
backward induction with DC-EGM in numba), \code{simulate.py} (monthly histories with common random numbers),
\code{panels.py} (the measurement layer of Section~\ref{sec:solution}), \code{moments.py} (the 97 targets with the data
scripts' estimators, the objective trim and the weights), \code{estimate.py} (the commands \code{evaluate},
\code{jacobian}, \code{polish}, \code{run}, \code{se}, \code{noise}), \code{borrowing\_wedge.py}, \code{bps\_check.py},
\code{hpc\_job\_a.py} and the staged checks in \code{tests/}. Each extension of 18 October is behind a flag of
\code{estimate.py} with its default set by the step-by-step evidence (\code{EXT\_DEFAULTS}: persistent CPS flows on,
response-error correction off, memory of the common shock off, borrowing wedge on, $\eta_{up}$ on). Outputs are
aggregate moments only; the PSID entry couples are git-ignored. The lifecycle model's estimates are not imported.

\paragraph{Compute and the HPC chain.} At the estimation size one evaluation takes 78 seconds on two threads (eight
problem-solves, 37{,}080 simulated couple-copies); at the full size about 510 seconds. The jobs run on the university
cluster in a separate clone of the repository with 32 cores, \code{--threads 2} (16 workers) and 64 GB, e-mail at begin,
end and failure, on a snapshot of the committed code (Decision~\ref{od:hpc}); the chain is in Table~\ref{tab:build}.

\begin{table}[htbp]\centering\small
\caption{The estimation chain on the cluster (\code{hpc\_job\_a.py})}\label{tab:build}
\begin{tabular}{L{3.6cm}L{11.0cm}}
\toprule
Job & What it does \\
\midrule
\code{Joint-test} & reproduces the objective at the start (11{,}733.7 over 85 moments, relative tolerance $10^{-3}$) and runs one Levenberg--Marquardt iteration \\
\code{Joint-recovery-gn} & pseudo-data recovery of all free parameters by Levenberg--Marquardt, start perturbed by 0.5 in logistic units \\
\code{Joint-gn-main} & Levenberg--Marquardt on the data from the closed-form start (15 iterations), Jacobian and sandwich standard errors at the end \\
\code{Joint-noise} & simulation variances of the moments at the candidate, five sets of common random numbers (the new weights) \\
\code{Joint-smm-main}, \code{-s2}, \code{-s3} & CMA-ES around the candidate with Jacobian-scaled initial steps, three seeds (3{,}500 evaluations each) \\
\code{Joint-gn-indep} & the independent-shock model, $\pi_c(\cdot) = 0$ and $m_H = 1$, by the same route \\
\code{Joint-se}, \code{Joint-jacobian}, \code{Joint-smm-block} & standard errors, Jacobian at the estimate, block weights as robustness \\
\bottomrule
\end{tabular}
\end{table}

% =====================================================================================================
\section{Decisions}\label{sec:decisions}

The PI took decisions D1--D18 on 18 October 2026, before the code, and D19--D24 on 18--19 October, after the first
evaluations of the coded model (\code{IGE/paperA\_model/results/}\allowbreak\code{diagnostics\_20261018.md}). Each states the choice and
the reason in one line; D25--D26 are options coded and switched off in the first estimate, to be decided at it. The body of the document implements them and refers to them by number.

\begin{enumerate}[label=\textbf{D\arabic*.}, ref=D\arabic*, leftmargin=*]
  \item \label{od:period} \textbf{Period.} Monthly, wife's ages 25--60 (420 months). \emph{Reason:} the joint loss is a
  same-month fact and the wife's response builds up over six months; the 480-month two-member solution of
  \citet{bacher2025joint} shows it is feasible.
  \item \label{od:prefs} \textbf{Preferences.} Common across races ($\phi_F$, $\phi_P$, $\theta_{05}$, $\theta_{617}$,
  $\phi_h$, $F$, $\kappa_S$); race enters through the hazards, the search friction, separation, inputs and wealth.
  \emph{Reason:} the joint loss is the same in both races, so racial differences in the response become a test of the
  environment.
  \item \label{od:child} \textbf{Child block.} Two slope scenarios (high: about $-9\%$ in adult earnings after a
  father's displacement, \citealp{oreopoulos2008intergenerational}; low: the smaller effects of
  \citealp{hilger2016parental}) and a term for both parents being present; the PSID associations are an upper bound.
  \emph{Reason:} the separation of the couple is a channel to the children, and the model translates parents'
  outcomes rather than estimating child development.
  \item \label{od:hpc} \textbf{HPC.} 32 cores, \code{--threads 2}, 64 GB (16 workers); 128 GB only with the PI's
  approval; never 64 cores; e-mail at begin, end and failure; job names \code{Joint-<what>}. \emph{Reason:} the PI's
  standing rules; the per-problem solve fits in about 4 GB per worker.
  \item \label{od:recall} \textbf{Job loss.} Permanent displacement only; the CPS targets use non-temporary job losses
  (excess 2.6 points on a base of 0.23\%). \emph{Reason:} temporary layoffs with recall are about a third of the CPS
  excess and are not the permanent loss the model studies.
  \item \label{od:trans} \textbf{Transitory wage shocks.} One draw per calendar year, held for twelve months, on two
  points per spouse; starting at BPS's variances and reduced by the part that is measurement error and other
  within-year sources, which enter only in the BPS measurement layer. \emph{Reason:} BPS's transitory shock is an
  annual object, monthly hourly wages are largely stable \citep{ganong2026earnings}, and a monthly i.i.d.\ shock would
  need a monthly standard deviation of 0.57 (husbands) and 0.39 (wives) (Section~\ref{sec:labor}).
  \item \label{od:wageproc} \textbf{Wage process.} Calibrated to BPS's Table 3 (permanent and transitory variances and
  covariances). \emph{Reason:} the reconciliation reproduces Table 3 in all 36 cells, and the BPS test then checks the
  mapping and the selection rather than refitting the wage process.
  \item \label{od:eras} \textbf{Policy schedules.} Three eras (1984--1996, 1997--2008, 2009--2021; 2022--23 with the
  last), each solved with its own stationary schedule, couples re-optimising unexpectedly at era boundaries.
  \emph{Reason:} the EITC, SNAP and UI changed over the PSID window; solving era by era keeps memory flat.
  \item \label{od:educ} \textbf{Education.} No education dimension at first; education enters through initial human
  capital and wealth; an education dimension is added only if the fit by education fails. \emph{Reason:} parsimony;
  four times the solves for a dimension the joint-loss moments do not need.
  \item \label{od:sep} \textbf{Separation.} The hazard is raised while the husband is out of work after a displacement.
  \emph{Reason:} it reproduces the separation response by race without an extra state.
  \item \label{od:credit} \textbf{Precautionary saving.} The discount factor fixed ($0.96$ a year), the borrowing limit
  estimated. \emph{Reason:} the same fixed preferences as the lifecycle model that shares this code; the wealth and
  consumption moments identify the limit.
  \item \label{od:cons} \textbf{Consumption targets.} PSID 1999--2021 only: the change after a single displacement (his,
  and hers) and the joint-minus-single difference are targets (its standard error, 0.073, is below the 0.08 rule); the
  BPS-panel $-0.138^{**}$ is validation only. \emph{Reason:} PSID expenditure starts in 1999, and the BPS-panel estimate
  depends on the sample.
  \item \label{od:husband} \textbf{Husband's search friction.} $\lambda_h(r)$ estimated jointly in the SMM; job finding by
  age validated against \citet{bacher2025joint}; one age slope added only if that validation fails. \emph{Reason:} the
  one-run rule and parsimony.
  \item \label{od:count} \textbf{Parameter count.} Thirty estimated parameters, with the cut list: a common $\delta_w$
  across races; $\phi_h$ calibrated to the husbands' part-time share; a common $\pi^{dn}$ for both spouses; $d_j$ if the
  re-employment earnings by spell length are fit by depreciation alone; the wife's lost-hours record in UI. The
  workplace match adds $\pi_c(2)$ (31 in Table~\ref{tab:params}), first on the cut list: merged with $\pi_c(1)$ if the
  Jacobian cannot separate them. \emph{Reason:} every parameter must move a named block of moments. (At the
  estimation size $s = 1, 2$ are merged, so $\pi_c(2)$ does not enter; with $\eta_{up}$ of Decision~\ref{od:eta}, 31
  parameters are free.)
  \item \label{od:welfare} \textbf{Welfare after separation.} The husband is not tracked after a separation.
  \emph{Reason:} the policy questions concern the household with the children; stated as a limitation of the welfare
  numbers.
  \item \label{od:search} \textbf{Search effort.} No separate search-effort state; ``looking'' is a measurement. The
  fallback, a non-employed state with and without search (one cost, one ratio of job finding), is used only if the base
  entry rate of non-employed wives (0.075 a month) and the SIPP build-up (0.019 to 0.054 over six months) cannot both be
  fit; decided at staged check (3). \emph{Reason:} parsimony; the acceptance margin produces the jump in looking.
  \item \label{od:awe} \textbf{Added-worker effect by age.} The effect of \citet{bacher2025joint} by age is untargeted
  validation (\code{IGE/50.CPSAddedWorkerByAge.py}, in progress). \emph{Reason:} no parameter depends on age, so the
  gradient is a test of the human capital, assets and horizon channels.
  \item \label{od:age} \textbf{Age range.} No extension to 64; the effect is validated at 26--35, 36--45, 46--55 and
  56--60. \emph{Reason:} the PSID copies are followed to the wife's 60, and the extension is not worth about 10\% more
  compute.
  \item \label{od:persist} \textbf{Persistent CPS flows.} The wives' entry and exit rates count a transition only if the
  old status held for two months before it and the new one for two months after it, in one rotation spell with the
  husband employed, identically in the data and the model; the response-error correction of
  \citet{poterba1986reporting} is dropped. \emph{Reason:} the model's monthly statuses contain no one-month reversals,
  and the correction implies a negative true entry rate at the data's employment level (Section~\ref{sec:targets}).
  \item \label{od:wedge} \textbf{Borrowing wedge.} Debt costs the saving rate plus 8.4 points a year, the G.19 real
  credit-card rate of 1995--2021 (11.4\%) less 3\%; a coded input, not estimated. \emph{Reason:} it lowers the
  objective at every borrowing limit tried and brings the young couples' debt share to the data
  (Section~\ref{sec:gov}); the 45--59 wealth shape stays an open misfit.
  \item \label{od:eta} \textbf{Accumulation by level.} $\pi^{up}_j\,(h/4)^{-\eta_{up}}$ with $\eta_{up}$ estimated
  (start 0.5). \emph{Reason:} the wage growth of the continuously employed falls with age faster than the bounded grid
  produces; the wage-growth block falls monotonically in $\eta_{up}$, though even 2.5 leaves most of the 45--54 gap
  (Section~\ref{sec:labor}).
  \item \label{od:ridge} \textbf{Couple type on the ridge; memory rejected.} $p_H$ and $m_H$ stay estimated, started on
  the CPS ridge at $p_H = 0.02$, $m_H = 15.95$; the memory of the common shock ($c_{mem}$) is kept as an option only.
  \emph{Reason:} on the ridge the PSID annual excess and the CPS leads and lags fit with no new parameter (joint-loss
  blocks 62 against 177 for the memory variant, whose excess one month around his loss is twice the data's).
  \item \label{od:hemp} \textbf{Husband's employment reported, not targeted.} The PSID employment rate of husbands by
  race (0.882 / 0.943) is in the fit table and not in the objective (model 0.999 / 1.000). \emph{Reason:} husbands leave
  work only by displacement and find a job in about four months, so no parameter reaches the target, which would only
  pull $\lambda_h$ down against the SIPP re-employment path (Section~\ref{sec:targets}).
  \item \label{od:sigma} \textbf{Taste-shock scale fixed.} $\sigma_\varepsilon = 0.05$, not estimated. \emph{Reason:}
  the exit rate stays at 0.2--0.35\% a month for scales from 0.05 to 0.30 against 2.6\%, while employment and entry rise
  away from the data (Section~\ref{sec:prefs}).
  \item \label{od:hersep} \textbf{Separation after her displacement: pending, off.} Coded as \code{--her-sep}: the
  separation hazard rises by $\sigma_{Jw}$ while she is entitled after her own displacement (\code{--her-sep-tie} sets
  it equal to $\sigma_J(r)$); target: still with the baseline husband in years 1--3 after her displacement, $-0.050$
  (s.e.\ 0.012), attrition-equalised (\code{IGE/64}, \code{IGE/65}, \code{IGE/66}). \emph{Status:} at the start values
  it moves the simulated response from $-0.001$ to $-0.019$ and the objective on the 89-moment set from 13{,}232.5 to
  13{,}162.8; decided at the first estimate on that set.
  \item \label{od:newjob} \textbf{Exit hazard in the first year of a job: pending, off.} Coded as \code{--new-job}: a
  state for a job started less than a year ago (left at rate 1/12) with an extra exit hazard $\xi_{nj}$; targets: SIPP
  monthly exit by tenure (0.021, 0.016, 0.0068 at 0--5, 6--11, 12+ months). \emph{Status:} at the start the model
  already over-predicts the first cell (0.035) and misses the third (0.0027); the option raises the first to 0.042 and
  the objective from 13{,}232.5 to 13{,}245.7, so the missing exits are those of established wives; decided at the first
  estimate (\code{results/extensions\_20261019.md}).
\end{enumerate}

\bibliographystyle{chicago}
\bibliography{model_document}

\end{document}
